\documentclass[10pt, a5paper,openany,twoside]{ksdp}

\usepackage{fontspec}
\usepackage{xunicode}
\usepackage{xltxtra}

\usepackage{graphicx}

\usepackage{geometry}
\geometry{paperwidth=13cm, paperheight=21cm}

\setlength{\baselineskip}{12pt}

%\usepackage[a5, frame]{crop} %mirko: schnittmarken, zentrieren

\usepackage{setspace} %Befehle: \singlespacing, \onehalfspacing, \doublespacing

\setlength{\topmargin}{-20,4mm} 
\setlength{\headheight}{9mm}
\setlength{\headsep}{5mm}
\setlength{\textheight}{\topskip+37\baselineskip} 

\setlength{\footskip}{0mm}

\setlength{\oddsidemargin}{-10,4mm}
\setlength{\evensidemargin}{-5mm}
\setlength{\textwidth}{94mm}
\setlength{\parindent}{5mm}

\pagestyle{fancy}
\usepackage{longtable}

\usepackage{csquotes}
\sloppy

\newcommand{\ksdpnumber}[1]{\fontspec{Georgia}\fontsize{8}{12}\addfontfeature{LetterSpace=5. 0}\selectfont #1\fontspec{Charis SIL}\fontsize{9}{12}\addfontfeature{LetterSpace=0.0}\selectfont} % mirko: anwendung: \ksdpnumber{gesperrter text} % auf schriftgrösse des folgetextes achten!

\newcommand{\ksdpbignumber}[1]{\fontspec{Georgia}\fontsize{9}{12}\addfontfeature{LetterSpace=5. 0}\selectfont #1\fontspec{Charis SIL}\fontsize{9}{12}\addfontfeature{LetterSpace=0.0}\selectfont}

\newcommand{\ksdpsmallnumber}[1]{\fontspec{Georgia}\fontsize{7}{12}\addfontfeature{LetterSpace=5. 0}\selectfont #1\addfontfeature{LetterSpace=0.0}\fontspec{Charis SIL}\fontsize{7}{12}\selectfont}

\newcommand{\ksdplargenumber}[1]{\fontspec{Georgia}\large\addfontfeature{LetterSpace=5. 0}\selectfont #1\addfontfeature{LetterSpace=0.0}\fontspec{Charis SIL}\large\selectfont}


\renewcommand{\chaptermark}[1]{%
 \markboth{#1}{}}

\setromanfont{Charis SIL}

\renewcommand{\contentsname}{} % mirko: überschrift des inhaltsverzeichnisses

\usepackage[NewCommands]{ragged2e}
%\usepackage[nohyph, russian]{polyglossia}
\usepackage{multirow}
\usepackage{wrapfig}
\usepackage{array}

\raggedbottom
\usepackage{multicol}
%%%%%%%%%%%%%%%%%%%%%%%%%%%%%%%%%%%%%%%%%%%%%%%%%%%%%%%%%%%%

\begin{document}

\fontspec{Charis SIL}
\fontsize{9}{12}\selectfont

% mirko: kopfzeile am kapitelanfang
\fancypagestyle{plain}{%
\fancyhf{} % clear all header and footer fields
\fancyfoot{}
\renewcommand{\headrulewidth}{0.0pt}
\fancyhead[LE]{\setlength{\RaggedRightParindent}{0cm}\raggedright\hspace{-10mm}\makebox[10mm][r]{\fontspec{Charis SIL}\ksdpbignumber{\thepage}}}
\fancyhead[RO]{\thepage}}

\thispagestyle{empty}

\title{Ла̄зэр ка̄ллса моаййнас}

%\hspace*{27.5mm}\parbox{6cm}{%
%\raggedright\vspace*{5\baselineskip}\vspace*{1mm}
%\fontspec{Profile}\fontsize{9}{12}\textbf{\textsc{kleine saamische schriften} — 1}\par
%\vspace*{.5\baselineskip}
%Herausgegeben von Michael Rießler\par
%und Elisabeth Scheller} 

\newpage
\thispagestyle{empty}

\vspace*{15mm}
\fontspec{Charis SIL}\fontsize{7}{12}\selectfont
\vspace*{8.5\baselineskip}
\parbox[t]{66mm}{

\begin{RaggedRight}

\noindent Lazar' Dmitrievič Jakovlev. Lāzer kāllsa moajjnas. Kɨr'jxa\\
lī Ēl'ce Nīna / Redaktion Nina Eliseevna Afanas'jeva und\\
Michael Rießler unter Mitarbeit von Elena Karvovskaja.\\
Berlin: Nordeuropa-Institut, \ksdpsmallnumber{2008}\\ 
Kleine saamische Schriften — \ksdpsmallnumber{1}\\[2mm]

\noindent Dieser Band entstand im Rahmen eines DoBeS-Projektes\\
zur Dokumentation der bedrohten kola\-saamischen Sprachen. Die DoBeS-Initiative wird von der \mbox{VolkswagenStiftung,} Hannover, gefördert. Weitere\\
Informationen zu DoBeS finden sich im Internet unter \mbox{http://www.mpi.nl/dobes.} \\[2mm]

\ksdpsmallnumber{1}. Auflage \ksdpsmallnumber{2008}

© Nordeuropa-Institut sowie die Autorinnen und Autoren\\[3mm]

Gestaltung — Anne-Christin Jyrch\\
Illustration — Dorothee Logen\\
Satz — Mirko Tietgen\\
Druck und Bindung — Druckhaus Berlin-Mitte\\[3mm]
Gesetzt in \XeLaTeX\ aus der Charis SIL\\[3mm]

\addfontfeature{LetterSpace=8.0}
ISSN\addfontfeature{LetterSpace=0.0}\ksdpsmallnumber{ 1866-5632}\\

\addfontfeature{LetterSpace=5.0}
ISBN\addfontfeature{LetterSpace=0.0}\ksdpsmallnumber{ 978-3-932406-91-1}
\addfontfeature{LetterSpace=0.0}

\end{RaggedRight}
}

\newpage
\thispagestyle{empty}

%%% schnipp --------

\vspace*{8\baselineskip}\vspace*{-0.9mm}
\vspace*{6pt}

\hspace*{11.6mm}\parbox{150mm}{
\begin{RaggedRight}
\fontsize{12}{18}\selectfont\textbf{\textit{\addfontfeature{LetterSpace=3.8}Л.\,Яковлев,\\
\fontsize{12}{24}\selectfont Н.\,Афанасьева, М.\,Рисслер (ред.)\addfontfeature{LetterSpace=0.0}}}

\fontsize{18}{24}\selectfont\textbf{\addfontfeature{LetterSpace=3}Ла̄зэр ка̄ллса моаййнас\addfontfeature{LetterSpace=0.0}}\vspace*{0.5mm}

\fontsize{12}{18}\selectfont\textbf{\addfontfeature{LetterSpace=3.8}Кырьйха лӣ Е̄льцэ Нӣна\addfontfeature{LetterSpace=0.0}}

\vspace*{14.5\baselineskip}
\vspace*{0.4mm}
%\fontspec{Profile}\fontsize{8}{12}\selectfont Nordeuropa-Institut der Humboldt-Universität | Berlin {2008}

\end{RaggedRight}
}

\fontspec{Charis SIL}\fontsize{9}{12}\selectfont

\newpage
\thispagestyle{empty}
\cleardoublepage

\tableofcontents
\fancyhead{}
\fancyhead[RO]{\thepage}
\fancyhead[LE]{}
\fancyfoot{}
\renewcommand{\headrulewidth}{0.0pt}

\cleardoublepage

\fancyhead{}
\fancyhead[RO]{\thepage}
\fancyhead[LE]{\hspace{-10mm}\makebox[10mm][r]{%
 \fontspec{Charis SIL}\ksdpbignumber{\thepage}}}
\renewcommand{\headrulewidth}{0.0pt}
\fancyfoot{}

\chapter*{\fontspec{Charis SIL}Яковлев Мыдтҍре Ла̄зэр\\
— \fontsize{12}{24}\textit{Нина Е.\,Афанасьева}}

\addcontentsline{toc}{chapter}{Яковлев Мыдтҍре Ла̄зэр\\
— \textit{Нина Е.\,Афанасьева}\vspace*{1.2mm}}

\fancyhead{}
\fancyhead[RO]{\thepage}
\fancyhead[LE]{\hspace{-10mm}\makebox[10mm][r]{\fontspec{Charis SIL}\ksdpbignumber{\thepage}}\hspace{0.5cm}\fontsize{9}{12}\selectfont{\fontspec{Charis SIL}Яковлев Мыдтҍре Ла̄зэр —} {\fontspec{Charis SIL}\addfontfeature{LetterSpace=5.0}\textit{Нина Е.\,Афанасьева}}\addfontfeature{LetterSpace=0.0}}
\fancyfoot{}
\renewcommand{\headrulewidth}{0.0pt}
\vspace{35.6mm}

\setlength{\RaggedRightParindent}{.5cm}\begin{RaggedRight}

\noindent Яковлев Мыдтҍре Ла̄зэр шэ̄ннтма лӣ шӯрр са̄мь пӣррсэсьт Кӣллт сыйтэсьт. 

Со̄нн чофта пуэраст тӣдӭ са̄мь кӣл, ко̄нн эннҍтэнҍ соннӭ аджесь, Вуафнэз Мыдтҍре Яковлев, Кӣлт сыйт са̄ммьля я е̄ннӭсь, Саваттӭй То̄ммьн Яковлева, туллэм са̄мь нызан. Сӣнэнҍ ля̄йй э̄ххтэмплоагкь па̄ррнэдтэ, пугк моаһтэнҍ пуэраст са̄мас соарнънэ. Школа райя па̄ррнэ тӣдтӭнҍ лышэ са̄мь кӣл. 

Ла̄зэр ля̄йй пуаррса, со̄нн о̄һпнувэ школасьт, маӈӈа о̄һпнувэ вуэпьсэй техникумэсьт Мурман ланӭсьт. 

Со̄нн робхужэ вуэппьсэенҍ оалк кла̄ссэнҍ Коардэгк сыйтэсьт, Луявьрэсьт. 
 
Куэссь со̄нн ля̄йй пенсиясьт, ро̄бхужэ со̄нн эфтэсьт мудтма соамегуэйм я
рӯшэгуэйм са̄мь кӣл а̄шэтҍ паяс пайнънэллмэнҍ, копче са̄мь са̄нӭтҍ Са̄мь-рӯшш соагкнэгк гуэйкэ. Е̄ннэ вӣгкэ лӣ аннтма тэнн лыһкэ. 

Ла̄зэр пуэраст тӣдӭ са̄мь моаййнсэтҍ, ко̄ппчма лӣ сӣнэтҍ я куадтма, штобэ мыйй вуајчемь пэ сӣнэтҍ олкэс лӯшшьтэ. 

Со̄нн лӣ ёаркса рӯшш кӣлэсьт кӣллт са̄мь кӣлле Аскольд Бажанов моаййнас »Вӣллькэсь пуаз«, ко̄нн кырьйха ля̄йй тетрадтӭ. Со̄нн ля̄йй авьтмусса, ке̄ тэнн тӯй че̄пельт ля̄йй пыййма са̄мь кӣлэсьт рӯшш килле.

\vspace*{\baselineskip}

\noindent Tэгкэ олкэс выййтма моаййнас лӣйенҍ кырьйха мунэнҍ, куэссь мыйй ро̄бхужэм Мыдтҍре Ла̄зэр ка̄ллсанҍ мӣн шӯрр са̄мь-рӯшш соагкнэгкь альн. Ла̄зэр мушштсэлэ тэйт са̄мь моаййнсэтҍ, мунн гыс сӣнэтҍ пыййе пымме э̄л – кырьйхэ са̄мас. Маӈӈа мунн о̄дзе тӣӈькэтҍ Ла̄зэр ка̄ллса моаййнсэтҍ олкэс лӯшшьтэм гуэйкэ. 

Tэгкэ моаййнас кырьйха ле̄в мӣн са̄мь алфавит я орфографья мӣлльтэ кӣлт са̄мь кӣлле. Мыдтҍре Ла̄зэр ка̄ллса пыййма лӣ е̄ннэ вӣгкэ я те̄дтэ мӣн са̄мь кӣлле аннҍтмэнҍ кырьйхэм на̄ль. Со̄н кырьйхэм я соарнънэм на̄лль кӣллт са̄мь кӣлэнҍ кӯдтъя лӣ тэйя моаййнсэнҍ, я пыннъюввма лӣ. 

Особливэ маӈӈьмусс со̄н моаййнас тэнн кырьйесьт »Ко̄ххт мунн оһпнувве« шигктэнне вузяһтҍ со̄н соарнънэм на̄ль. Tэдта лӣ адтҍ ныдтҍ шэ моаййнас мӣнэ гуэйкэ, нэ со̄нн лӣ э̄ллм моаййнас Мыдтҍре Ла̄зэр баяс, со̄н о̄һпнуввэм ыге баяс. Tэнн моаййнас кырьйхэ ля̄йй Г.\,М.\,Керт \ksdpnumber{1960} ыгесьт Луявьрэсьт фонетическэ транскрипция мӣлльтэ. Я ныдтҍ шэ со̄нн пыййма лӣ тэнн моаййнас рӯшш кӣлле ӣжесь кырьесьт »Са̄м кӣл соарнънэм на̄ль« (Образцы саамской речи). 

Са̄мь кӣл те̄дт-олма, Мэ̄һкал Рӣсслер, лыгкэ я кырьеһтҍ тэнн моаййнас кырьйе са̄мь-рӯшш соагкнэгкъ лыстэтҍ я лӯшьтэ со̄н олкэс. 

\end{RaggedRight}\setlength{\RaggedRightParindent}{0cm}

\vspace*{5.1mm}


\begin{FlushLeft}
\hspace*{-3mm}
\begin{tabular}{r b{5cm}}
\multirow{8}{44mm}{\includegraphics[width=4.6055cm]{foto_300dpi_sw}}
	& \\
	& \\
	& \\
	& \\
	& \\
	& \\
	& \\
	& \\
	& \\
	& \\
	& \\
	& \textit{Лазарь Дмитриевич Яковлев},\\
	& \ksdpnumber{1918–1993}\\
\end{tabular}
\end{FlushLeft}


\chapter*{\fontspec{Charis SIL}Иммель а̄йя}
\addcontentsline{toc}{chapter}{Иммель а̄йя}

\fancyhead{}
\fancyhead[RO]{\thepage}
\fancyhead[LE]{\hspace{-10mm}\makebox[10mm][r]{\fontspec{Charis SIL}\ksdpbignumber{\thepage}}\hspace{0.5cm}\tiny{\fontspec{Charis SIL} }\hspace{0.25cm}\fontsize{9}{12}\selectfont{}}
\fancyfoot{}
\renewcommand{\headrulewidth}{0.0pt}
 \vspace*{44mm}

\setlength{\RaggedRightParindent}{.5cm}\begin{RaggedRight}

\noindent Ко̄ххт Иммель энҍтэ ло̄нҍтэтҍ нюнӭтҍ я ю̄лькэтҍ. Чуэннь пӯдӭ, Иммьлесьт кэ̄же мо̄джесь нюнҍ я ю̄лькэтҍ. Иммель, анҍт мыннӭ рӯппьсэсь руэчкэдҍ. Ӣджь чуэннь лӣ се̄рэ, чоаһпесь нюнӭнҍ, рӯппьсэсь руэчкэгуэйм. Чуэннь лӣ мо̄джесь лоаннҍт. 

Маӈӈа пӯдӭ нюххч, Иммьлесьт кэ̄же ю̄лькэтҍ я нюнҍ. Иммель энҍтэ нюххча чоаһпесь ю̄лькэтҍ, чоаһпесь нюнҍ, оанҍхэсь пэдтша, роавас суетҍ я вӣллькэсь понцэтҍ. Нюххч лӣ шӯрр вӣллькэсь лоаннҍт, со̄нэсьт палл ча̄дзь-чуххч. Нюххч куэннт не̄лльй маннӭ. 

То̄ххт пӯдӭ Иммьлесьт кэ̄дже мо̄джесь кыррьй понцэтҍ, чогк нюнҍ я рӯппьсэсь ю̄лькэтҍ, гӯ чуэнньгэсьт. Иммель са̄ррн: »То̄ххт, мунн тоннӭ анта мо̄джесь кыррьй понцэтҍ, чоаһпесь чогк пасьтлэсь нюнҍ, а рӯппьсэсь ю̄льк еввла.« То̄ххт се̄ннтэль: »Мугка чоаһпесь ю̄лькэтҍ эмм ва̄льт,« – я кыррьтэль ю̄лькэха. Иммель кыҏтхэмушшэ о̄ййкэй кудтҍ пя̄ла пэдтша чоаһпесь ю̄лькэтҍ. Адтҍ то̄ххт эйй ма̄тҍ ва̄ннҍцэ е̄ммьнэ мӣлльтэ, сонэсьт ю̄льк ле̄в маӈӈень. То̄ххт пуэраст вуайй ча̄зесьт я кугкь нуэһк ча̄зь вульн. То̄ххт куэннт эфт манҍ, кӯһт манҍ, ко̄ллм манҍ.

Югке ыгка со̄нн луэшшт ча̄дза эфт манҍ пуфч гуэйкэ. Куэссь то̄ххт куэннт эфт манҍ, со̄нэсьт эйй лӣннче то̄ххтэ – алка. Куэссь куэннт кӯһт манҍ, со̄нэсьт куэдтай э̄ххт алкэнч. Куэссь куэннт ко̄ллм манҍ, со̄нн элкадҍ кӯһт манҍ. Танна со̄нэсьт ле̄в кӯһт то̄ххтэ – алка. То̄ххт лыһк пе̄ссь я̄ввра, я̄вьр рыннта. Коашшьк ыгенҍ со̄нн куэннт рыннт поавьн э̄л, пе̄ссь гурре со̄нн пуадт ло̄сстэнне. Вуэйй рынт гурре, нюнэнҍ чогкаст е̄ммьнэ я роаӈькэсь кя̄сс алльтла я нюнэнҍ чогкаст эвтэсьт, ныдтҍ со̄нн пе̄зант пе̄ссь райя. Моайнас пугк. 

\end{RaggedRight}\setlength{\RaggedRightParindent}{0cm}

%\vspace*{3.5cm}
%\hspace*{-2.5cm}
\vspace*{2\baselineskip}
\begin{center}
\includegraphics{600_rabe_2st_fliegend}
\end{center}

\chapter*{\fontspec{Charis SIL}Ко̄ххт о̄зэ рӣммьн ва̄ресьт савехь мӯр}
\addcontentsline{toc}{chapter}{Ко̄ххт о̄зэ рӣммьн ва̄ресьт савехь мӯр}

\fancyhead{}
\fancyhead[RO]{\thepage}
\fancyhead[LE]{\hspace{-10mm}\makebox[10mm][r]{\fontspec{Charis SIL}\ksdpbignumber{\thepage}}\hspace{0.5cm}\fontsize{9}{12}\selectfont{\fontspec{Charis SIL}Ко̄ххт о̄зэ рӣммьн ва̄ресьт савехь мӯр}}
\fancyfoot{}
\renewcommand{\headrulewidth}{0.0pt}
\vspace*{34.6mm}

\setlength{\RaggedRightParindent}{.5cm}\begin{RaggedRight}

\noindent Рӣммьн о̄зэ ва̄ресьт савехь мӯр. Ё̄дӭ ва̄рь мӣлльтэ я ло̄гэ: »Касьт лӣ ва̄ресьт шӣһтлэс мӯрр савьхэ гуэйкэ?« Вуэйвэсь паянт паяс я аппс: ка̄ррнас – алка апснэнҍ. Пырр мӯр ва̄ннц я са̄ррн: »Тала лӣ шӣг савехь мӯрр. Пырр ва̄рь воаннҍцэ мугка мӯр эмм коаввнма. Тэдта мӯрр лӣ самэ шӣг мӯрр, мунн со̄н чӯпа.« Ка̄рьнэсь – а̄ка пайенҍ куӆшкӯдӭ, штэ рӣммьн та̄һт тэнн мӯр чӯһпэ. Ка̄рьнэсь – а̄ка са̄ррн соннӭ: »Рӣммьн, ель чӯп тэнн мӯр, мунэстҍ алка ле̄в.« Рӣммьн гыс са̄ррн соннӭ: »Анҍт мыннӭ эфт алка, танна мӯр эмм чӯп.« Ка̄рьнэсь – а̄ка сӯһпей пайенҍ вэлас эфт алка. Рӣммьн воальтэ тэнн я по̄рэ. Е̄ннэ лӣ ва̄нас а̄ййкэ мэ̄нэ, рӣммьн аввта по̄ҏсувэшкӯдӭ, пӯдӭ тэнн шэ мӯрре, пырс ёдтал я са̄ррн: »Пугк весь ва̄рь пырс кэрре, никасьт мугка мӯр эмм коаввнма. Адтҍ тэнн мӯр чӯпа.« 

Ка̄рьнэсь – а̄ка эйй кӣҏҏшэнч я са̄ррн: »Рӣммьн, ель чӯп тэнн мӯр. Ка̄ррнас парнайт ма̄ййкак.« Ка̄рьнэсь е̄ннҍ вуэллькэдҍ вэлас вя̄л эфт алка. Рӣммьн тэнн гыс по̄рэ я ёадтэй ӣжясь тӯе мӣлльтэ. Ё̄дӭ, ё̄дӭ, по̄ҏсувэшкӯдӭ я аввта мушштэй тэнн мӯр баяс. Пӯдӭ со̄нн тэнн шэ мӯрре, вылшхалл паяс я са̄ррн: »Адтҍ чӯпа тэнн мӯр, пэря эмм коаввнма никасьт.« 
 Ка̄рьнэсь – а̄ка са̄ррн: »Рӣммьн, ель чӯп тэнн мӯр. Мунн вя̄л эфт алка анта.« Ка̄рьнэсь – а̄ка вуэллькэдҍ вя̄л эфт алка вэлас. Рӣммьн по̄рэ я уйтэ. Ка̄рьнэсь кырьтэсьт. Ка̄рьнэсь ря̄гк: »Кроӈьк, кроӈьк, кроӈьк. Ке̄ лӣ пуэдтлэнч?« 

Рӣммьн пуэдтэльт, мунэсьт ко̄ллм алка по̄рэ. Коале мунн эмм леһче аннтэнч алкатҍ, со̄нн ля̄һч пэ тэнн мӯр те̄ҏҏпэ. Ка̄рьнэсь са̄ррн агкнесь: »Манҍтэ ля̄к то̄нн то̄ӆӆкхэм, а̄ка. Рӣмьнэсьт элля ни оакшэ, ни ныййп. Со̄нн эйй вуэй е̄дтэ тэнн мӯр. »Ка̄рьнэсь кыррьтэль, э а̄ка кӯдтэй мӯр альн оаррэ. Рӣммьн пӯдӭ, вуагэсьт, есскэль я паяс вылшхалл. Тамьпэ кӯһт ка̄ррнас – алка кӯдтъенҍ. 

Адтҍ мунн тэнн мӯр чӯпа. Ва̄рь мӣлльтэ кэ̄рре, кэ̄рре, никасьт эмм коаввнмэнч мугка шувв мӯр. 

А Ка̄рьнэсь – а̄ка са̄ррн: »Рӣммьн, то̄нн эгк вуэй чӯһпэ тэнн мӯр, то̄нэсьт элля ни оакшэ, ни ныййп.« Рӣммьн юррьтэль, юррьтэль я вуэмэшьт, штэ ке̄-ля̄ннч ля̄йй пуэдтлэнч я ма̄тҍхэнч ка̄ррнас агка, штэ со̄нэсьтҍ элля ни оакшэ, ни ныййп. Юррьтэль, юррьтэль, то̄ӆӆькъенч ля̄йй, штэ ка̄рьнэсь ля̄йй пуэдтлэнч я ка̄ррнас – агка ма̄тҍхэнч. 

Рӣммьн судтэ ка̄ррнас э̄л я са̄ррн: »Адтҍ мунн со̄н по̄ра.« Вуагэсьт, вуагэсьт я ламьп ря̄ввна кэ̄рркэдҍ, гудэль ка̄ллма. Се̄йп пруткшэдтӭ я кыбршувэ я ле̄жант. Ка̄рьнэсь кырьтэ, кырьтэ, кроӈькэсьт: »Кроӈьк, кроӈьк.« Уйнэсьт: рӣммьн лӣ ка̄ллмэнч ла̄мьп ре̄вьнэсьт. Юррт: »Адтҍ мунн по̄ра со̄н.« 

Рӣммьн ля̄шш, ня̄лльм лӣ а̄ввнэнч, ню̄ххчемь ыдтъя, чильмэтҍ лӣ оамьпхэ, гудэль ёаммэнч лӣ. Ка̄рьнэсь кырьтэ, кроӈькэсьт: »Кроӈьк, кроӈьк.« Ся̄йвэ, рӣмьн гоаррэ ва̄ннц, юррт: »Адтҍ ню̄ххчемь со̄нэсьт по̄ра!« Э̄ллтэнҍт, лышэ ню̄ххчемь тувэ нуэһкэ, рӣммьн кадцесьт со̄н ча̄пхэсьт. Те̄ннҍтэльт со̄н я са̄ррн: »Мунн то̄н по̄ра.« 

Ка̄рьнэсь ря̄гк: »Рӣммьн, ель по̄р мун, ель поадэһтҍ понцэтҍ.« Ка̄рьнэсь вульн юррт: »Алка пе̄знэ!« Рӣммьн, ель по̄р мун тэста. Юрьтэшкуэдэтҍ, штэ рӣммьн лӣ по̄ррэнч ка̄ррнас тэста. Мунн ӣджь ёадтъя роадтэ, э то̄нн гыс мӣлльтэ ка̄ль. 

Рӣммьн юррт, зоабэль, мэ̄нн мунн алка со̄н ке̄ссе? Ло̄сса лӣ ке̄ссе. Луэшта со̄н, со̄нн ӣджь манн тоагэ, мунн тамьпэ е по̄ра со̄н. Ня̄льмэдэсь а̄ввьдэсьтэ, ка̄рьнэсь нюччкэй я кыррьтэль. Кыррт, кроӈькаст: »Кроӈьк, кроӈьк.« 

Рӣммьн кя̄һч, ка̄рьнэсь вӯййлэдтӭ со̄н. Ныдтҍ рӣммьн эйй вуайя ка̄ррнас по̄ррэ, ка̄рьнэсь со̄н шӣгктэнне вӯййлэдтӭ. Моайнас пугк. 

\end{RaggedRight}\setlength{\RaggedRightParindent}{0cm}
 
\chapter*{\fontspec{Charis SIL}Ко̄ххт кӯль эллькэнҍ\\
ча̄зесьт е̄лле}
\addcontentsline{toc}{chapter}{Ко̄ххт кӯль эллькэнҍ ча̄зесьт е̄лле}

\fancyhead{}
\fancyhead[RO]{\thepage}
\fancyhead[LE]{\hspace{-10mm}\makebox[10mm][r]{\fontspec{Charis SIL}\ksdpbignumber{\thepage}}\hspace{0.5cm}\fontsize{9}{12}\selectfont{}}
\fancyfoot{}
\renewcommand{\headrulewidth}{0.0pt}
\vspace*{35.6mm}

\setlength{\RaggedRightParindent}{.5cm}\begin{RaggedRight}

\noindent Кӯль пӯдтӭнҍ Иммьле, кэ̄джьлэнҍ, касьт Иммель поаһк сыйетҍ е̄лле? Авьтмусса пӯдӭ Иммьле лӯсс я са̄ррн: »Иммель, лӯшьт мун ме̄рр ча̄зесьт я соавьй ча̄зесьт е̄лле. Ме̄рэсьт мунн алка шэ̄ннтэ, э соавьй ча̄дза алка мӣйнэтҍ ко̄дтэ. Соавьй ча̄зесьт аллькэв шэ̄ннтэ лӯз – алка. Маӈӈа сыйй кӣррсэв соавьй ё̅гэ мӣлльтэ ме̄рре. Ме̄рэсьт сыйй шэ̄ннтэв шӯрр кӯль райя. Иммель лыгкэ ныдтҍ, ко̄ххт кэ̄же лӯсс. Лӯсс лӣ мо̄джесь тоавас ме̄рр кӯлль, чӯввесь чӯмэгуэйм, рӯппьсэсь вӯйй цӯпэнҍ. Соавьй ё̅гэнҍ лӯсс пыннь ӣжесь алкатҍ, тэнн гуэйкэ со̄нн поабпхалл вӯӈӈькэ. 

Вуэсск лӣ лӯз вӣлльй. Вуэсск я̄лл соавьй ча̄зесьт ыгегуэйм, вуэсск лӣ мо̄джесь кӯлль, чӯм ле̄в руэнн – чоаһпесь чуаягуэйм. Иммель энҍтэ соннӭ мугка кыррьй коавь. Вуэсскан лӣ па̄ннҍ я чӯллья̄вьр кӯлль. Кыдта вуэсск коадт сӯйн лӯһтэ сызнҍ. Ныдтҍ Иммель ся̄дэ вӯсскнэ е̄лле. Моайнас пугк. 

\end{RaggedRight}\setlength{\RaggedRightParindent}{0cm}

\chapter*{\fontspec{Charis SIL}Моайнас рӣмьн\\
я нюэммель баяс}
\addcontentsline{toc}{chapter}{Моайнас рӣмьн я нюэммель баяс}
\fancyhead{}
\fancyhead[RO]{\thepage}
\fancyhead[LE]{\hspace{-10mm}\makebox[10mm][r]{\fontspec{Charis SIL}\ksdpbignumber{\thepage}}\hspace{0.5cm}\fontsize{9}{12}\selectfont{}}
\fancyfoot{}
\renewcommand{\headrulewidth}{0.0pt}
\vspace*{35.6mm}

\setlength{\RaggedRightParindent}{.5cm}\begin{RaggedRight}

\noindent Рӣммьн я нюэммель коавншуввенҍ ва̄ресьт. Нюэммель са̄ррн: »Рӣммьн, аллькэпь эфтэсьт чӯһкэ мӣлльтэ ёадтӭ.« Рӣммьн вуагк, вуагк я ескант, маӈас кя̄һч, куллталл. Нюэммель каджь: »Рӣммьн, мэййт то̄нн ескнак я куллтэлак?« Рӣммьн са̄ррн: »Мунэсьт пэдэсьт еввла чальм« – я сӣйнънэшкӯдӭ. Нюэммель юррт: мунн ныдтҍ шэ тӣшшладта со̄н. Нюэммель карр чӯһкэ мӣлльтэ раст-раст чӯһкэ. Тэнн пя̄лла чӯһка карр, ескант, куллталл. Нымьп пя̄лла чӯһка карр, куллталл. Рӣммьн са̄ррн: »Мэййт то̄нн, нюэммель, раст-раст чӯһка карак я куллтэлак?« Нюэммель са̄ррн: »Чӯһка вӯлленҍ элля чуэкас« – я сӣйнънэшкӯдӭ. Мунн ныдтҍ шэ рӣмьн тӣшшлэһтӭ. Кӯһта сӣйнэнҍ. Нюэммель сӣйнэнҍт, штэ рӣмьнэсьт пэдэсьт еввла чальм, э рӣммьн сӣйнэнҍт, штэ чӯһка вӯлленҍ элля чуэкас, тэнн гуэйкэ карр нюэммель раст-раст чӯһка. 

\end{RaggedRight}\setlength{\RaggedRightParindent}{0cm}

%\hspace*{-2.5cm}
\vspace*{2\baselineskip}
\begin{center}
\includegraphics{600_hase_1st}
\end{center}

\chapter*{\fontspec{Charis SIL}Моайнас рӣмьн я ко̄рк баяс}
\addcontentsline{toc}{chapter}{Моайнас рӣмьн я ко̄рк баяс}

\fancyhead{}
\fancyhead[RO]{\thepage}
\fancyhead[LE]{\hspace{-10mm}\makebox[10mm][r]{\fontspec{Charis SIL}\ksdpbignumber{\thepage}}\hspace{0.5cm}\fontsize{9}{12}\selectfont{}}
\fancyfoot{}
\renewcommand{\headrulewidth}{0.0pt}
 \vspace*{44mm}
 
 \setlength{\RaggedRightParindent}{.5cm}\begin{RaggedRight}

\noindent Рӣммьн я ко̄ррк вуэсстлэннҍтэнҍ шӯрр ла̄мьп альн. Ко̄ррк ля̄йй се̄ннтэнч рӣмьн э̄л. Рӣммьн ко̄ркэсьт я мудта ло̄нҍтэнҍ по̄рэнҍт мэнӭтҍ, удць алкатҍ. Ко̄ррк тэнн гуэйкэ тувэ рӣмьн вуэпьсэ. Ко̄ррк я са̄ррн: »Рӣммьн, та̄так то̄нн мун се̄льк альн кырьтсе? Мунн то̄н ма̄тҍха кыррьтэ.« Рӣммьн са̄ррн: »Мунн тула та̄һтэ кыррьтэ то̄н ка̄вьр альн.« Ко̄ррк са̄ррн; »Нюам мыннӭ се̄льк э̄л, кӯһт ке̄бпель пый пырр ча̄пэх, кӯһт маӈӈь ке̄бплэтҍ пый кудтҍ пя̄лла пэдташ, э се̄ййп анҍ пруткшуввант маӈем пе̄льт.« Рӣммьн оаштэ ыштэ ко̄рк ка̄вьр э̄л, ко̄ххт ко̄ррк па̄кэ со̄ннӭ ышштэ. Ко̄ррк кыррьтэль, пэ̄йенҍт вӯссьт ва̄рь э̄л, поакэ рӣммьнэ нюччкъе. Рӣммьн кэдче вэлас, вусьт поаввьнэ, оарр эйй лыгкэнҍт. Ко̄ррк кырьтэ пырр рӣмьн, ныммпэшьт кырьтэ пырр рӣмьн. Рӣммьн паенҍт вуэйвэсь, ко̄ррк каджь со̄нэсьт: »На, ко̄ххт мэ̄нэ кыррьтҍмушш?« Рӣммьн цӣӆькэ: »Кыррьтҍмушш мэ̄нэ эйй шӣгктэнне.« Ко̄ррк каджь: »То̄нн сӣйпэнҍ ёртътэкь?« Рӣммьн гыс са̄ррн: »Эмм кӣрркнэнч юррьтле, ко̄ххт вэлас пӯдтӭ, вусьт поаввна кэ̄ррктэ. 

Адтҍ кыррьтлак, ель вуаййладтҍ сӣйпэнҍ ёртътэ. Рӣммьн о̄дтӭсьт нюамэ ко̄ррка ка̄вьр э̄л. Маӈӈҍ ю̄лькэтҍ ныдтҍ шэ пые, ко̄ххт е авьтма вуэр, сӣйпэнҍ кӣһчлэнтэ ёртътасьтэ. Ко̄ррк кыррьтэль, паяс пэ̄йенҍт, мӣлкнесьт пэ̄йенҍт шӯрр ва̄рь э̄л. Я со̄нн поакэ рӣммьнэ кыҏтхэ. Рӣммьн нюччкэй я вӯлькэ вэлас, гудэль ке̄дҍк. Вуаййлэдтӭ сӣйпэнҍ ёртътэ, кэ̄рркэдҍ вусьт кокорнэһка, ю̄лльк мэ̄нэ раст. Рӣммьн вя̄йкэрт, пруткшувэ я ля̄шш, эйй лыгкэнҍт нимӣ. Ко̄ррк кырьтэ пырр рӣмьн. Рӣммьн эйй лыгкэнҍт, гӯ ёамма лӣ. Ко̄ррк вуэллгэнне кырьтэ пырр рӣмьн, кэ̄джель со̄нэсьт: »Е̄лльмэнҍ ля̄к?« Рӣммьн паенҍт вуэйв, ня̄льм а̄венҍт, эйй вуайя нимэ̄нн це̄ӆӆькэ. Ко̄ррк ся̄йвэ рӣмьн лӯз я мӣлкнесьт цӣӆькэ: »Ама, алкак вя̄л ло̄нҍтэнҍ мэнӭтҍ по̄ррэ я лоанҍт алкатҍ суэллнэ?« Тэль ко̄ррк кыррьтэль ӣжесь ла̄мьп гоаррэ, касьт ля̄йй пе̄ссь сайй мэнӭгуэйм. Моайнас пугк. 

\end{RaggedRight}\setlength{\RaggedRightParindent}{0cm}

\enlargethispage{15\baselineskip} 
\vspace*{2\baselineskip}
%\hspace*{6cm}
\begin{center}
\includegraphics{600_rabe_1st}
\par
%\vspace*{-11\baselineskip}
%\hspace*{-1.5cm}
%\includegraphics{600_rabe_1st}
\end{center}


\chapter*{\fontspec{Charis SIL}Ко̄ххт Иммель энҍтэ\\
ноавьтэтҍ чӯрьвэтҍ я пӣльетҍ}
\addcontentsline{toc}{chapter}{Ко̄ххт Иммель энҍтэ ноавьтэтҍ чӯрьвэтҍ\\
я пӣльетҍ\vspace*{1.2mm}}

\fancyhead{}
\fancyhead[RO]{\thepage}
\fancyhead[LE]{\hspace{-10mm}\makebox[10mm][r]{\fontspec{Charis SIL}\ksdpbignumber{\thepage}}\hspace{0.5cm}\fontsize{9}{12}\selectfont{}}
\fancyfoot{}
\renewcommand{\headrulewidth}{0.0pt}
\vspace*{35.6mm}

\setlength{\RaggedRightParindent}{.5cm}\begin{RaggedRight}

\noindent Сэ̄ррв пӯдӭ Иммьле я са̄ррн: »Анҍт мыннӭ шӯрр чӯрьвэтҍ ко̄ммьтэсь лоабпэгуэйм я са̄рркмэгуэйм.« Иммель энҍтэ соннӭ шӯрр мо̄джесь чӯрьвэтҍ. Сэ̄ррв кэ̄же Иммьлесьт кугкесь пӣльетҍ, штобэ пуэраст куллэ. Иммель энҍтэ соннӭ кугкесь пӣльетҍ. 

Пӯдзъе Иммель энҍтэ эйй шӯрр пӣльетҍ я мо̄джесь са̄ррк чӯрьвэтҍ я оанҍхэсь пӣдчна, гӯ сэрвэсьт. Иммель энҍтэ пӯдзъе югке – налшэм сэксмэтҍ: вӣллькэсь, чоаһпесь, руччкесь, кыррьй. 

Нюэммель пӯдӭ, кэ̄же Иммьлесьт шӯрр чӯрьвэть, манҍтэ сэрвэсьт ле̄в. Иммель энҍтэ соннӭ шӯрр чӯрьвэтҍ, я нюэммьля каррэль сӣнэтҍ кӣһчлэнҍт. Карр, вусьт мӯрэ пашшнант, ёрэ шӯрр чӯрьвэгуэйм, по̄нншадт, пэ̄ййнэ эйй вуэй. Иммель уййн, штэ нюэммель эйй вуэјт сӣнэгуэйм шӣӆкхуввэ я са̄ррн: »Нюэммель, эйй быдтӭ тоннӭ шӯрр чӯрьвэтҍ. Мунн тоннӭ анта кугкесь кулльесь пӣльетҍ. Тэгк лӣннченҍ пэрямп, мэ̄нн шӯрр чуэрьв.« Ныдтҍ нюэммель элькэ е̄лле кугкесь пӣльегуэйм, оанҍхэсь пӣджьнэнҍ, удлэсс ю̄лькэгуэйм, вӣллькэсь вӯгкесь сэксмэнҍ та̄лльва я ке̄ссь па̄ль се̄ра сэксмэнҍ. Моайнас пугк. 

\end{RaggedRight}\setlength{\RaggedRightParindent}{0cm}

\chapter*{\fontspec{Charis SIL}Вуэррев я ме̄һцьлэнч}
\addcontentsline{toc}{chapter}{Вуэррев я ме̄һцьлэнч}

\fancyhead{}
\fancyhead[RO]{\thepage}
\fancyhead[LE]{\hspace{-10mm}\makebox[10mm][r]{\fontspec{Charis SIL}\ksdpbignumber{\thepage}}}\fancyfoot{}
\renewcommand{\headrulewidth}{0.0pt}
 \vspace*{44mm}
 
 \setlength{\RaggedRightParindent}{.5cm}\begin{RaggedRight}

\noindent Ме̄һцьлэнч вӯлькэ мя̄һца пе̄ннганҍ. Пе̄нна коавнэ вуэррев пе̄ссь сай, пе̄ссь сызьн лӣйенҍ вуэррев – алка я вуэррев – е̄ннҍ. Пе̄ннэ элькэ угкэ. Вуррев – е̄ннҍ са̄ррн алкатҍ: »Тыйй оарре пе̄сь сызьн мӣлкнесьт, елле кӯввьлсалле!« Эфт алкне цӣӆьке: »То̄нн мӣлкнесьт луэштэдак вэлас, маӈӈьмусс вуэкьс э̄л. Кӣч че̄д пе̄зь куэһцьве, манҍтэ лӣ ме̄һцьла я пе̄ннэ.« Вуэррев алкэнч пӯдӭ пе̄дзь ла̄ддва, касьт ля̄йй пе̄ссь я са̄ррн: »Ме̄һцьла лӣ ва̄птэгк о̄ллмэнч, пиннчак лӣ о̄дт, пысс пай вуаӆӆк, пе̄нна лӣ пӯйтэсь.« Е̄ннӭсь са̄ррн: »Тыйй, па̄ррнэ, елле палле, оарре мӣлкнесьт. Анҍ ме̄һцьла пе̄зь куалхталл. Со̄нн вэххтэнне уййт.« 

Вуэррев алка ыййдэнҍ, ӣнца аввта пӯдӭ ме̄һцьла, элькэ пе̄зь куалхтэллэ. Пе̄ннэ угк. Е̄ннӭсь са̄ррн: »Луэштэд вэлас, выльшэһтҍ, манҍтэ ме̄һцьла лӣ пуадтма.« Вуэррев – алкэнч нӣе вэлас маӈӈьмусс вуэкьс э̄л я кя̄һч: о̄ллмэнч лӣ вуэммь ка̄хтнэсьт, рӯсст пысс, вуэммь савехь, пе̄ннэ лӣ не̄лльк. Нюччк вусьт пе̄дзе, ыррашт, угк ӣннҍеннӭ. Вуэррев – алка карэ паяс пе̄зь е̄ннӭсь лӯз я са̄ррн: »Ме̄һцьленч лӣ ва̄йваш, ка̄ххтан лӣ вуэммь, рӯсст пысс, вуэммь савехь. Пе̄ннэ лӣ не̄лльк.« 

Е̄ннӭсь са̄ррн: »Тыйй, па̄ррнэ, оарре, елле лыгкнэгке. Тэдта ме̄һцьла мӣнэтҍ коаввн.« Ӣджь вӯлькэ о̄дзэ чӯһкас мӯрэ мӣлльтэ, коаӈк ле̄в пе̄зь, кӯз сӯгктэ я алямп. Вуэррев – е̄ннҍ пӯдӭ моазэ я са̄ррн па̄рнэдӭсь: »Адтҍ аллькэпь мӣлкнесьт нюччкэ кя̄джьлэнне мӯрэсьт мӯрра. Мунн эвтас алка нюччкэ, э тыйй мӣлльтэ. Ныдтҍ сыйй нючкэшкудтэнҍ: вуэррев – е̄ннҍ эвтэсьт нюччк я маӈас кӣчант, ко̄ххт пуэдтӭв мӣлльтэ со̄н па̄ррнэ. Ныдтҍ вуэррев – е̄ннҍ воаеһтҍ е̄ллем сай. Адтҍ вуэррев – \enlargethispage{15\baselineskip} алка пӯдтӭнҍ о̄дт сайя. А ме̄һцьлэнч лышэ куалхталл курас пе̄зь, а пе̄ннэ угк курас пя̄дзэ. Ныдтҍ пынньчесс вуэррев пӣсьтэ ма̄дэсьт ӣжесь алкатҍ. Моайнас пугк. 

\end{RaggedRight}\setlength{\RaggedRightParindent}{0cm}

\vspace*{2\baselineskip}
%\hspace*{1.5cm}
\begin{center}
\includegraphics{600_eichhorn_1st}
\end{center}

\newpage

\vspace*{6\baselineskip}
%\hspace*{-2.5cm}
\begin{center}
\includegraphics{600_hecht_1st}\hspace*{0.5cm}\includegraphics{600_hecht_2st}
\end{center}
%\par
%\vspace*{6\baselineskip}
%\hspace*{2.5cm}
%\includegraphics{hecht_2st}

\chapter*{\fontspec{Charis SIL}Ныгкешь баяс}
\addcontentsline{toc}{chapter}{Ныгкешь баяс}

\fancyhead{}
\fancyhead[RO]{\thepage}
\fancyhead[LE]{\hspace{-10mm}\makebox[10mm][r]{\fontspec{Charis SIL}\ksdpbignumber{\thepage}}\hspace{0.5cm}\fontsize{9}{12}\selectfont{}}
\fancyfoot{}
\renewcommand{\headrulewidth}{0.0pt}
 \vspace*{44mm}
 
 \setlength{\RaggedRightParindent}{.5cm}\begin{RaggedRight}

\noindent Иммель энҍтэ нэ̄һкша кугкесь роаӈьк, руэнн чӯмэтҍ, чоаһпесь чуаягуэйм, штобэ эйй кӯсстъече пэ чӣллк ча̄зесьт я ра̄зе сызнҍ. Ныгкешь вуэйвэсьт ле̄в пассьтлесь па̄нҍ, тэдта лӣ паннай кӯлль. Пугк кӯль ӣлленҍ ча̄зесьт, э ныгкешь эйй та̄һтэнч ча̄зесьт е̄лле. Иммель о̄ййкэй со̄н ча̄дза кӯрьвэгуэйм. 

Адтҍ ныгкешь цӯпэсьт ле̄в се̄ӈькэсь сӯррь тахта, гудэль куэрьв я нэ̄ммдуввэв »куэррьв та̄хьт«. Вя̄л ныгкешь вуэйвэсьт лӣ »ка̄ла-та̄ххьт«. Нызан, кыппьтмэнҍ нэ̄һкаш вуэйв, пыйнэв ка̄лагоанҍцэтҍ по̄ррэ. Нӯрр мо̄ллда нимэшкӯдӭ нэ̄һкаш вуэйв я нӣлэсьт тэнн та̄хьт, ко̄нн адтҍ коадчетҍ \hspace*{-0.2mm}»ка̄ла-та̄ххьт«. Кӯль эллькэнҍ Иммьле е̄ввэ, штэ ныгкешь сӣнэтҍ вайяд, штэ ныгкешь поарряй лӣ, пукэтҍ поарр. Е̄на пукэ ваййдувэнҍт нэ̄һкшэсьт вуэсск. Танна Иммель энҍтэ вуэсскнэ чо̄гк пассьтлэсь сэ̄дткэтҍ ка̄вьр альн. Ныгкешь со̄н эйй вуэјт нӣллэ. Тэйе сэ̄дткэгуэйм вуэсск пя̄сстай нэ̄һкшэсьт. Вя̄л Иммель ся̄дэ, штэ ма̄н ке̄жесьт пӯдтҍмэнҍ нэ̄һкшэсьт па̄нҍ кэчнэв. Ныдтҍ кӯль пе̄зсэлнэв нэ̄һкшэсьт. Э̄ввтэль пугк кӯль соарнънэнҍ. А суэввель е̄ннэ ке̄ллсэлэнҍт кӯлетҍ Иммель баяс. Танна вуэсск вӯесьт суэввьле я лыгкэ пассьтлэсь сэ̄дткэнҍ райк суэввель ню̄ххчмэсьт. 

Адтҍ мыйй суэввель по̄ррмэнҍ уййнэпь суэввель ню̄ххчмэсьт тэнн райк, ко̄нн лыгкэ соннӭ вуэсскан, штобэ суэввель никуэссегенҍ эйй ке̄ллсэлче пэ. Моайнас пугк. 

\end{RaggedRight}\setlength{\RaggedRightParindent}{0cm}

\chapter*{\fontspec{Charis SIL}Моайнас кӣмьн баяс}
\addcontentsline{toc}{chapter}{Моайнас кӣмьн баяс}

\fancyhead{}
\fancyhead[RO]{\thepage}
\fancyhead[LE]{\hspace{-10mm}\makebox[10mm][r]{\fontspec{Charis SIL}\ksdpbignumber{\thepage}}\hspace{0.5cm}\fontsize{9}{12}\selectfont{\fontspec{Charis SIL}Моайнас кӣмьн баяс}}
\fancyfoot{}
\renewcommand{\headrulewidth}{0.0pt}
 \vspace*{44mm}
 
 \setlength{\RaggedRightParindent}{.5cm}\begin{RaggedRight}

\noindent Ля̄йй эфт о̄ллмасьт шӯрр кӣммьн. Ка̄ллесь, ко̄ххт ёадтсэлэшкуадт мя̄һца, кӯль шылле, мӯрэ гоаррэ ва̄ррэ – пуадт ке̄ммна я каджь со̄нэсьт: »Ко̄ххт пэря лыһкэ, мӣ вуайй ва̄ресьт вуасста пуэдтӭ? Манҍтэ пӣдтҍ шаннт ва̄ресьт?« Ка̄ллсэсьт ля̄йй пӣрас: а̄ка я мо̄джесь нӣййт. Ӣлленҍ, ӣлленҍ сыйй я нӣййт элькэ кэ̄джьтэдтэ нымьп сыййта, касьт лӣ руэнньса е̄лле. Эйй по̄р, нюэзельт вуадт, пай олкэс кя̄һч. Аджесь са̄ррн со̄нэнҍ; »Мэ̄ййт, нӣйтням, эгк ля̄ ве̄ссель? Мэ̄нн баяс юртак?« Нӣййт са̄ррн эженӭсь: »Мунн шэ̄ннҍтэ шӯрр, адтҍ та̄та кӯйе выйтэ. А касьт мунн коавна па̄рьн, ко̄нн мунн шоабпша?« Аджесь са̄ррн: »То̄нн кӣҏш вя̄л эйй кугкь а̄ййке.« Э ӣджь юррт: »Мунн кӣмьнэнҍ, шӯрр кӣмьнэнҍ соарнсаста.« 

Е̄кьнэнч пӯдӭ, аджесь вӯлькэ кӣмьнэнҍ соарнънэ: »Кӣммьн, то̄нн, кӣммьн, то̄нн мун югке вуэр куллтэллэкь я ве̄кьхэллекь мыннӭ. Адтҍ мунэсьт лӣ шӯрр пӣдтҍ.« Кӣммьн каджь: »Мӣ тоннӭ шэ̄нтэ?« – Мыннӭ эйй шэ̄ннтэнч нимӣ. Мун нӣййт пэчксэлант, кэпшсэлант, элля руэняс соннӭ ва̄ррь моазэсьт е̄лле. Поаллтлэнне мӣнэгуэйм шӯрр я̄вьр тӯгкенҍ лӣ сыййт, касьт нӯрраш ве̄ссьлэнне е̄ллев. Эгк вуайче ко то̄нн тэнн сыйт ыгка ке̄ссе мӣн е̄ллем сайе? Кӣммьн са̄ррн: »Тэдта элля ло̄ссесь лыһк. Мунн тэнн то̄н кэ̄джьмуж ӣллькъя.« 

Куэдӭсьт вуадтъенҍ, ля̄йй цӣгк ыйй. Кӣммьн чо̄ннкӯдӭ пэйель куэӆьц я мэ̄нэ тэнн сыййта, касьт нӯрраш сӣрр пэр, тэсьт нюччкенҍ я сӣрренҍ. Кӣммьн мӣлкнесьт ке̄сськӯдӭ тэнн пэр, т, касьт нӯраш сӣрэ, тэнн гоаррэ, касьт нӣййта эйй ле̄ннч руэняс е̄лле. 

Па̄ррнэ нюччкенҍ, сӣрренҍ, олкэс выййтэнҍ: ев тӣдҍ, касьт сыййт ля̄йй. Куэррьев чуэннчев тэйт шэ сэенҍ, нӯтҍ рынтэсьт оаллкэнҍ лоабпэв, тэсьт са̄ййм ле̄в коадзэнч олк альн. Эйй кӯһкень ка̄ллса пэҏт луннҍ оарр. 

\enlargethispage{5\baselineskip}
Ка̄ллса ӣнца цоаввн нӣйтэсь: »Адтҍ то̄нн мэ̄н, томтэслувв сӣрьегуэйм. Нӣййт пӯдӭ, суаӈэ сӣррэм пэ̄ҏьтэ. Па̄ррнэ я нӣйт вэлльенҍ соннӭ вуэсста, тӣрвхэдтэв, со̄нэсьт нэ̄м кэ̄джев. Ныдтҍ, ка̄ллса, агка нӣййт томтэслувэ тэйе е̄лльегуэйм, койт кӣммьн кӣзе ыгка тэнн пэҏт луз. Маӈӈа нӣййт коавнэ ӣджяс мо̄джесь па̄рьн, я сыйй ноајјтлэдтэнҍ. Е̄лешкудтэнҍ руэнньсэнне. Сӣнэнҍ шэ̄ннҍтэнҍ па̄ррнэ. О̄дт сыйтэсьт эллькэнҍ е̄лле. Пукэнҍ ля̄йй лыһк, пугк е̄лешкудтэнҍ пуэраст. 

А ка̄ллесь пӯдӭ ке̄ммна, па̄ссьпужэ я вя̄л пэрямп элькэ пынне кӣмьн. Моайнас пугк. 

\end{RaggedRight}\setlength{\RaggedRightParindent}{0cm}

\vspace*{2\baselineskip}
%\hspace*{1cm}
\begin{center}
\includegraphics{600_kessel_2st}
\end{center}

\chapter*{\fontspec{Charis SIL}Мэ̄нн гуэйкэ оадзэ о̄ллмэгуэйм эйй е̄ль}
\addcontentsline{toc}{chapter}{\protect Мэ̄нн гуэйкэ оадзэ о̄ллмэгуэйм эйй е̄ль}

\fancyhead{}
\fancyhead[RO]{\thepage}
\fancyhead[LE]{\hspace{-10mm}\makebox[10mm][r]{\fontspec{Charis SIL}\ksdpbignumber{\thepage}}}
\fancyfoot{}
\renewcommand{\headrulewidth}{0.0pt}
\vspace*{35.6mm} 

\setlength{\RaggedRightParindent}{.5cm}\begin{RaggedRight}

\noindent Эйй лӣннче руэняс оадзкэ э̄ххтэ ва̄ресьт е̄лле. Нике̄ со̄нэнҍ эйй оаҏ-тэдтмэнч, никуэссь со̄нн нике̄з эйй ря̄ххтэдтма я эйй лӣйе ке̄з нюэссь тӯетҍ лыһкэ. Со̄нн та̄тэ поагэ ва̄ресьт уййтэ, штобэ о̄ллмэгуэйм вӯййксэлнэ. 

Пугк та̄тэ ӣжесь-налла лыһкэ, югке сайе ӣжесь оадзэ ке̄бплэтҍ лӯшшьтэ. Ла̄мьпэсьт ӣжесь ка̄нҍц кӣзе, ӣжесь луз лыгктэдэ. Шӯрр ва̄птэгэтҍ оашшьтэй соннӭ, штобэ ка̄ннҍцэсь пугквесь уйнэч я куллтлэнч. Юррьтэль ре̄һкь я сыннӭсь оадзэ, ко̄ххт пукэтҍ э̄ллмантэсьт пе̄ссьтьэ. 

О̄ллма я̄ллмушш чофта эйй мыльса соннӭ, я пукэтҍ оадзэтҍ се̄вьнэсь ва̄ресьт коадче ӣжесь гурре. Эфт о̄ллма ма̄йкэ, нымьп выеһтҍ, куаллмнэ тя̄һтэдэ, ня̄льянт э̄л па̄нӭтҍ я кэ̄нцэтҍ та̄һкад. Тэнн шэ сыйтэсьт вуэммь нуэййт ӣле. 

Нуэййтъюшшэ со̄нн эйй ма̄һтэнч, пай вӣресь бубен мӣлльтэ пляссэй, тэнн гуэйкэ, штэ ӣжесь бубен эйй лӣйе. Оадзэ воальтэ со̄н ӣжесь кӣдэ вуэла. Энҍтэ мугка па̄һкмуж, штобэ вӯййкэсьт соннӭ лушшэ, э оаз со̄н па̄ссьпаһт. Оаз ма̄тӭ е̄ннэ нуэййтъюшшэ я та̄тэ вӣресь кӣдэгуэйм ня̄льянт о̄ллма ма̄ййкэ, кӯ ля̄йй роаввса пӯдзэ чуэррьвэ я та̄л па̄ннӭ. Копче оаз пугк ӣжесь я̄кхэллъетҍ, эллькэнҍ юррьтэ, ко̄ххт нюэссь о̄ллмасьт сыйе пе̄сстэдтэ. Юррьтэнҍ, юррьтэнҍ я мугка па̄һкмуж лыһкенҍ, штэ тэдта олмэнч нюэзельт лыһк я эйй ма̄тҍ оадзэ бубен вӯлленҍ пляссъе. 

А нуэййт па̄һкмуж вуэла ӣжесь рыст ця̄кэсьт. Тэль сӣнэнҍ нимӣ эйй выййтма. Тэнна о̄ллма пӯррь са̄м тӯльйсэллэнҍ: со̄н эйй ве̄ррьвуд вӯййкъенҍ. Нуэййт ӣжесь рыст гуэйкэ вуэсстланнт, э оаз кӯдтэй туэресь пя̄лла. Ноэ пусьтэльт оаз роаммшэдэ, штэ ке̄һпьсэнне сӣнэнҍ пе̄сстэдэ. 

О̄ллма ев аллкэнч ӣнха оаррэ, ев аллка со̄н нюэссь соагэтҍ кӯннҍтэ. Пӯдтӭнҍ я цӣӆӆькэнҍ оадзкэ: »Мыйй тоннӭ ва̄лт энҍтэмь я мыйй шэ то̄нэсьт поагэ ва̄лльтэпь.« Я выйхэнҍ оадзэ моазэ ва̄рра е̄лле. Моайнас пугк. 

\end{RaggedRight}\setlength{\RaggedRightParindent}{0cm}

\vspace*{2\baselineskip}
%\hspace*{1cm}
\begin{center}
\includegraphics{600_kessel_1st}
\end{center}

\chapter*{\fontspec{Charis SIL}Ко̄ххт мунн о̄һпнувве}
\addcontentsline{toc}{chapter}{\protect Ко̄ххт мунн о̄һпнувве}

\fancyhead{}
\fancyhead[RO]{\thepage}
\fancyhead[LE]{\hspace{-10mm}\makebox[10mm][r]{\ksdpbignumber{\thepage}}\hspace{0.5cm}\fontsize{9}{12}\selectfont{\fontspec{Charis SIL}Ко̄ххт мунн о̄һпнувве}}
\fancyfoot{}
\renewcommand{\headrulewidth}{0.0pt}
 \vspace*{44mm}
 
 \setlength{\RaggedRightParindent}{.5cm}\begin{RaggedRight}

\noindent Э̄ввтэль мӣн аджь я е̄ннҍ соарнэнҍ, штэ мӣн сыйтэсьт эйй лӣййма шко̄ла са̄мь парна гуэйкэ. Мӣн сыйтэсьт, Кӣллт са̄мь сыйтэсьт шко̄ла ля̄йй уже маӈӈа тӯййшэнч. Тэнн шко̄ласьт о̄һпнуввенҍ са̄мь парна кӯһтэмплоагкь нийтнедтӭ я выдтэмплоагкь па̄ррьшэнче, сӣнэгуэйм мунн тэсьт шэ о̄һпнувве. 

Адтҍ мушта, штэ, куэссь мунн о̄һпнувве пе̄рвэ кла̄ссэсьт, о̄һпнувве пуэраст. Мушта, куэссь ёаме Ылльй Ладэмыр Ленин, мыйй учтельницэнҍ воаррдэм ва̄рра, пыдтӭмь кӯсс куэһцьвэтҍ, штобэ соннӭ лыһкэ венок. 

Маӈӈа мунн о̄һпнувэшкӯдтӭ эвтас, не̄лльй кла̄сс мунн лӯпьтэ Кӣллт сыйтэсьт. Куэссь мунн о̄һпнувве тамьпэ, мунн во̄дтэдтӭ Куэллнэгке. Тамьпэ мэ̄нэ пионерэтҍ коппчнэдтмушш, я мунн лӣйе гӯ делегаһт. Тамьпэ мунн соарнэ, ко̄ххт о̄һпнуввэв са̄мь парна Кӣллт сыйтэсьт, маӈӈа пӯдтӭ моазэ, эллькэ эвтас о̄һпнуввэ. Куэссь мунн лӯппьтэ не̄лльй ыгь, мунн юрьтэшкӯдтӭ: эвтас о̄һпнуввэ, эле вӯллькэ пӯдзэтҍ пынне. 

Пе̄рвэ ыгь мун поаһкэшь пынне ӣжян пӯдзэтҍ. Мун де̄дӓ, Вуафнэз Ша̄шка, куэссь со̄нн ё̅дӭ рай мун, танна со̄нн цӣӆӆьке, штэ адтҍ быдт о̄һпнуввэ. Я мӣн соаменҍ е̄ввла вуэппьсуввма о̄ллма, темэтҍ быдт о̄һпнуввэ. 

Мунн юррьтэ, юррьтэ я эллькэ эвтэс о̄һпнуввэ. 

Маӈӈа пуэдтэль мӣн сыййта Осипов Эвван, Вуэннҍтре алльк, тэдта ля̄йй пе̄рвэ са̄мь вуэппьсэй. Со̄нн соарнэ, штэ Мӯрмнэсьт лӣ вуэппьсэем пэ̄ҏҏт, касьт алькэтҍ вуэпьсэ са̄мь па̄ррнатҍ, я мунн вӯллькэ тэнн шко̄лайе. Тамьпэ мунн о̄һпнувве подготовительнэ отделениясьт. Кӯһт ыгь мунн о̄һпнувве. Пе̄рвэ ыгь мунн о̄һпнувве я лӯппьтэ выдант я кудант кла̄сс, нымьп ыгь мунн лӯппьтэ кыджянт кла̄сс. 

Куэссь мунн о̄һпнувве Мӯрмнэсьт, мӣнэтҍ, са̄мь па̄ррнэтҍ, вӯӆкхэшь сыйтэ мӣлльтэ вуэпьсэ шӯрр о̄ллматҍ. Тамьпэ мӣнэнҍ лӣйенҍ са̄мь кнӣга, букваррь, кнӣга ло̄гкэм гуэйкэ ля̄йй, я лӣйенҍ мудта кнӣга са̄мь кӣлле лыһкма. Мунн танна вуэпьсэ Коардэгэсьт, Коардэгк сыйтэсьт. Куэссь мунн мэ̄ннӭ тэнн сруэк, пӯдтӭ маӈӈа моаст Мӯрмнэ, уже эллькэ эвтас о̄һпнуввэ. 

Куэссь мунн лӯппьтэ Мӯрманскэ педагогическэ училища, мунн эллькэ ро̄бхушшэ ӣжян сыйтэсьт. Пе̄рвэ ма̄тҍхэ Коардэгэсьт, маӈӈа Чудзья̄вьрэсьт, адтҍ вуэпьса Луя̄вьрэсьт. 

Танна ля̄йй нюэзе о̄һпнуввэ, мэ̄нн адтҍ. Адтҍ мӣн советскэ правительства я коммунистическэ па̄ртия пя̄дцлушшэв мӣнэ баяс. Мыйй, удць са̄мь народӭнч, ныдтҍ шэ я̄ӆчемь пуэраст я лӣһчемь кырьйхэенҍ. 

Адтҍ тыста, Луя̄вьрэсьт, са̄мь парна о̄һпнуввэв шӯрр шко̄ласьт, кӯ лӣ ке̄дҍкэсьт лыһкмэнч. Тэдта шко̄ла лӣ ко̄ллм этажэнҍ. Адтҍ Луя̄вьрэсьт тӯййшувв са̄мь я ыжэм парна гуэйкэ шко̄ла – интернаһт, касьт аллькэв сыйй вуэпьсэ, ко̄ххт е̄лле. Тэсьт парна аллькэв мэ̄ннэ производственнэ пра̄ктикум, ле̄в лыһкма мастерскэ. Адтҍ са̄мь парна я ыжэм парна я рӯшш, когк тэсьт, Луя̄вьрэсьт ле̄в, аллькэв ро̄бхушшэ е̄ммьнэ альн, аллькэв куаййвэ е̄ммьнэ, ко̄ххт се̄йе се̄мятҍ, я о̄һпнуввэ, ко̄ххт пынне пӯдзэтҍ, штобэ ванҍса коадхэ, штобэ пӯдзэ шэнтченҍ пэрямп, штобэ колхосс о̄ллма вя̄л пэрямп е̄лешкуэдтӭв. 

Тэнн пайнэмушшэ ве̄кяһт советскэ правительства я коммунистическэ па̄ртия. Сыйй лыһкэв пай пэрямп, штобэ са̄мь о̄ллма я̄ӆченҍ пэрямп.

\vspace*{\baselineskip}
 
\noindent(Кырьйха лӣ тэнн Мыдтҍре Ла̄зэр моаййнас Г.\,М.\,Керт \ksdpnumber{1960} ыгесьт Луя̄вьрэсьт.)

\end{RaggedRight}\setlength{\RaggedRightParindent}{0cm}

\enlargethispage{5\baselineskip}
\vspace*{2\baselineskip}
\begin{center}
\includegraphics{600_hase_2st}
\end{center}

\clearpage
\hbox{}
\thispagestyle{plain}

\chapter*{\fontspec{Charis SIL}Удць соакнэгка\\
— \fontsize{12}{24}\textit{Михаэль Рисслер}}
\addcontentsline{toc}{chapter}{\protect Удць соакнэгка\\
— \textit{Михаэль Рисслер}}

\fancyhead{}
\fancyhead[RO]{\thepage}
\fancyhead[LE]{\hspace{-10mm}\makebox[10mm][r]{\fontspec{Charis SIL}\ksdpbignumber{\thepage}}\hspace{0.5cm}\fontsize{9}{12}%\selectfont{\fontspec{Charis SIL}Михаэль Рисслер}
}
\fancyfoot{}
\renewcommand{\headrulewidth}{0.0pt}
\vspace*{35.6mm} 

\setlength{\RaggedRightParindent}{.5cm}\begin{FlushLeft}

\noindent Ниже приводится список всех саамских слов, а также словоизменительных и словообразовательных морфем, которые встречаются в тексте »Ко̄ххт мунн о̄һпнувве«.

Русский перевод текста опубликован в сборнике саамских текстов под редакцией Г.\,М.\,Керт »Образцы саамской речи. Материалы по языку и фольклору саамов Кольского полуострова (кильдинский и иоканьгский диалекты)« (Москва-Ленинград: Наука, \ksdpnumber{1961}). На странице \ksdpnumber{25} сборника приводится тот же самый текст в фонетической транскрипции.

Перевод словоформ в данном списке слов довольно краток, часто он дан лишь в виде грамматических объяснений. Таким образом, приведенный список слов не заменяет собой словарь, он может лишь помочь при самостоятельном изучении саамского текста.

В списке слов приводятся не только основные формы слова, но и его алломорфы (например окончание \textbf{-aһта}, которое является алломорфом окончания $\rightarrow$\,\textbf{-ха} абесс. ед.\,ч.). Русский перевод дается только для одного вхождения слова или морфемы, для всех остальных форм того же самого корня, словоизменительного или словообразовательного аффикса (алломорфов) указывается только соответствующая начальная форма. Указывается также часть речи (например \textit{глаг.}) и чередования в основе \mbox{|-уэлк-, -уэллк-, -ӯльк-|} для глагола \textbf{вӯллькэ} ›идти‹). Для некоторых исторически сложившихся словоформ помимо перевода дается исходная морфема (например, для порядкового числительного \textbf{выдант} \mbox{›пятый‹ $\leftarrow$\,выдт, -ант).} Ниже приводятся используемые сокращения.

Более полный перевод большинства встречающихся в текстесловоформ, а также примеры предложений вы можете найти в словаре Р.\,Д.\,Куруч, Н.\,Е.\,Афанасьева и Е.\,И.\,Мечкина »Са̄мь-рӯшш соагкнэһкь« (Москва: Русский язык, \ksdpnumber{1985}). Часть этого словаря представляет собой краткое грамматическое описание наиболее важных парадигм. При самостоятельном изучении грамматики саамского языка могут быть также полезны списки парадигм в словаре »Удць сáмиса̄мь са̄мь-сáми соагнэгка« Пекка Саммаллахти, Анастасия Хворостухина (Ohcejohka: Girjegiisá, \ksdpnumber{1991}). Большинство слов из списка вы также можете найти в словаре Г.\,М.\,Керт »Словарь саамско-русский и русско-саамский« (Ленинград: Просвещение, \ksdpnumber{1986}).\par

\end{FlushLeft}\setlength{\RaggedRightParindent}{0cm}

\begin{multicols}{2}

\setcounter{unbalance}{16}

\raggedcolumns % dieser sinnlos anmutende befehl muss hier stehenbleiben, weil multicol nur registerhaltig setzt, wenn er hier steht.

\RaggedRight

\textbf{-а} \textit{оконч.} \ksdpnumber{1}л. ед.\,ч. наст.\,вр. 

\textbf{-а} \textit{оконч.} илл. ед.\,ч. 

\textbf{-а} \textit{суфф.} ласк.\,ф. 

\textbf{аджь} \textit{сущ.} \mbox{|ажь-|} отец

\textbf{адтҍ} \textit{нареч.} теперь, в настоящее время

\textbf{-aһта} \textit{оконч.} \mbox{$\rightarrow$\,-ха}

\textbf{-ак} \textit{оконч.} \ksdpnumber{2}л. ед.\,ч. наст.\,вр. 

\textbf{алльк} \textit{сущ.} \mbox{|альк-, эльк-, аллк-|} сын

\textbf{аллькэ} \textit{глаг.} \mbox{|алк-, аллк-, -элльк-|} начинать

\textbf{аллькэв} \textit{глаг.} \mbox{$\rightarrow$\,аллькэ, -эв}

\textbf{алькэтҍ} \textit{глаг.} \mbox{$\rightarrow$\,аллькэ, -этҍ}

\textbf{альн} \textit{послел.} на, над, из, с

\vspace*{\baselineskip}

\textbf{баяс} \textit{послел.} о, об, про

\textbf{букваррь} \textit{сущ.} \mbox{|-арь-, -арр-|} букварь

\textbf{быдт} \textit{глаг.} (\ksdpnumber{3}л. ед.\,ч.) нужно, надо, необходимо

\textbf{-бэдтӭ} \textit{оконч.} \ksdpnumber{2}л. мн.\,ч. наст.\,вр. 

\textbf{ванҍса} \textit{нареч.} меньше \mbox{($\leftarrow$\,ва̄нас, -а)}

\textbf{ва̄рра} \textit{сущ.} \mbox{$\rightarrow$\,ва̄ррь, -а}

\textbf{ва̄ррь} \textit{сущ.} \mbox{|-а̄рь-, -а̄рр-|} лес

\textbf{ве̄кьхэ} \textit{глаг.} \mbox{|ве̄кьха, ве̄кяһт|} помогать, помочь \mbox{($\leftarrow$\,ве̄һкь, -х-, -э)}

\textbf{ве̄кяһт} \textit{глаг.} \mbox{$\rightarrow$\,ве̄кьхэ}

\textbf{воаррдэмь} \textit{глаг.} \mbox{$\rightarrow$\,воаррьдэ, -эмь}

\textbf{воаррьдэ} \textit{глаг.} \mbox{|воаррьда, воарряд|} беречь, стеречь, охранать

\textbf{во̄дтэдтӭ} \textit{глаг.} \mbox{|-эд-, -эдт-, -эдҍ-|} бывать, иметь доступ \mbox{($\leftarrow$\,во̄дт-, -дт-, -э)}

\textbf{Вуафнэсс} \textit{собст.} \mbox{|-эз-, -сс-|} Вуафнэсс, Афанасий

\textbf{вӯлкхэ} \textit{глаг.} \mbox{|вӯлкха, вӯлкаһт|} отправить, послать, направить \mbox{($\leftarrow$\,вӯлльк-, -х-, -э)}

\textbf{вӯлкхэшь} \textit{глаг.} \mbox{$\rightarrow$\,вӯлкхэ, -эшь}

\textbf{вӯллькэ} \textit{глаг.} \mbox{|-уэлк-, -уэллк-|} идти, пойти, направиться

\textbf{Вуэннҍтре} \textit{собст.} \mbox{|о.\,ч.|} Вуэннҍтре, Андрей

\textbf{вуэппьсуввма} \textit{глаг.} \mbox{$\rightarrow$\,вуэппьсуввэ, -ма}

\textbf{вуэппьсуввэ} \textit{глаг.} \mbox{|-ув-, -увв-, -уввь-|} воспитываться; заниматься \mbox{($\leftarrow$\,вуэпьс-, -увв-, -э)}

\textbf{вуэппьсэй} \textit{сущ.} воспитатель \mbox{($\leftarrow$\,вуэпьс-, -эй)} 

\textbf{вуэппьсэем} \textit{прилаг.} \mbox{$\rightarrow$\,вуэппьсэй, -эм}

\textbf{вуэпьса} \textit{глаг.} \mbox{$\rightarrow$\,вуэпьсэ, -а}

\textbf{вуэпьсэ} \textit{глаг.} \mbox{|вуэпьса, вуэпяст|} воспитывать, воспитать

\textbf{выдант} \textit{числ.} пятый \mbox{($\leftarrow$\,выдт, -ант)}

\textbf{выдтэмплоагкь} \textit{числ.} пятнадцать \mbox{($\leftarrow$\,выдт, -эмп-, лоагкь)}

\textbf{вя̄л} \textit{нареч.} ещё

\vspace*{\baselineskip}

\textbf{гӯ} \textit{союз} как

\textbf{гуэйкэ} \textit{послел.} для, за, из-за

\vspace*{\baselineskip}

\textbf{делегаһт} \textit{сущ.} \mbox{|-ат-, -һтҍ-|} делегат

\vspace*{\baselineskip}

\textbf{-е} \textit{оконч.} \ksdpnumber{3}л. мн.\,ч. прош.вр.

\textbf{-е} \textit{оконч.} иллат. ед.\,ч. 

\textbf{-е} \textit{оконч.} \mbox{$\rightarrow$\,-э}

\textbf{ёадтӭ} \textit{глаг.} \mbox{|ёад-, ёадт-, ё̄д-|} ехать

\textbf{ёаме} \textit{глаг.} \mbox{$\rightarrow$\,я̄мме}

\textbf{е̄ввла} \textit{глаг.} они нет \mbox{($\leftarrow$\,е̄в ля̄й)}

\textbf{ё̅дӭ} \textit{глаг.} \mbox{$\rightarrow$\,ёадтэ}

\textbf{е̄лешкуэдтӭв} \textit{глаг.} \mbox{$\rightarrow$\,е̄лле, -шкуэдтҍ-, -эв}

\textbf{е̄лле} \textit{глаг.} \mbox{|я̄л-, я̄лл-, йӣл-|} жить, доживать, дожить

\textbf{е̄ммьнэ} \textit{сущ.} \mbox{|о.\,ч.|} земля, почва

\textbf{-енҍ} \textit{оконч.} \mbox{$\rightarrow$\,-энҍ}

\textbf{-енҍ} \textit{оконч.} \ksdpnumber{3}л. мн.\,ч. прош.вр.

\textbf{е̄ннҍ} \textit{сущ.} \mbox{|е̄нҍ-, я̄ннҍ-, я̄нн-, е̄ннҍ-|} мать

\textbf{-енч} \textit{суфф.} \mbox{$\rightarrow$\,-энч}

\textbf{-есьт} \textit{оконч.} \mbox{$\rightarrow$\,-эсьт}

\textbf{-етҍ} \textit{оконч.} \mbox{$\rightarrow$\,-этҍ}

\textbf{-еха} \textit{оконч.} \mbox{$\rightarrow$\,-эха}

\vspace*{\baselineskip}

\textbf{ӣжян} \textit{местоим.} себя (акк.\,/\,ген. ед.\,ч.) \mbox{($\leftarrow$\,ӣджь, -ан)} 

\textbf{интернаһт} \textit{сущ.} \mbox{|-ат-, -һтҍ-|} интернат

\vspace*{\baselineskip}

\textbf{-йе} \textit{оконч.} илл. ед.\,ч.

\vspace*{\baselineskip}

\textbf{касьт} \textit{местоим.} где; откуда (лок. ед.\,ч.)

\textbf{ке̄ддҍк} \textit{сущ.} \mbox{|-е̄дҍк-, -е̄ддк-|} камень

\textbf{ке̄дҍкэсьт} \textit{сущ.} \mbox{$\rightarrow$\,ке̄ддҍк, -эсьт}

\textbf{кӣлл} \textit{сущ.} \mbox{|-ӣл-, -ӣлль-|} речь, язык

\textbf{кӣлле} \textit{сущ.} \mbox{$\rightarrow$\,кӣлл, -е}

\textbf{Кӣллт} \textit{собст.} \mbox{|-ӣлт-, ӣлльт-|} Кӣллт, кильдинский

\textbf{кла̄сс} \textit{сущ.} \mbox{|о.\,ч.|} кла̄сс

\textbf{кла̄ссэсьт} \textit{сущ.} \mbox{$\rightarrow$\,кла̄сс, -эсьт}

\textbf{кнӣга} \textit{сущ.} \mbox{|о.\,ч.|} книга

\textbf{коадхэ} \textit{глаг.} потерять, утерять

\textbf{Коардэгк} \textit{собст.} \mbox{|-эг-, гкь-|} Коардэгк, Воронье

\textbf{Коардэгэсьт} \textit{собст.} \mbox{$\rightarrow$\,Коардэгк, -эсьт}

\textbf{когк} \textit{местоим.} (ном. мн.\,ч.) которые, какие

\textbf{ко̄ллм} \textit{числ.} \mbox{|-о̄лм-, -о̄льм-|} три

\textbf{колхосс} \textit{сущ.} \mbox{|-оз-, -ссь-|} колхоз

\textbf{коммунистическэ} \textit{прилаг.} коммунистический

\textbf{коппчнэдтмушш} \textit{сущ.} \mbox{|-уж-,-шшь-|} сбор \mbox{($\leftarrow$\,коппч-, -нэдт-, -мушш)}

\textbf{ко̄ххт} \textit{нареч.} как

\textbf{кӯ} \textit{местоим.} который, какой (ном. ед.\,ч.)

\textbf{куаййвэ} \textit{глаг.} \mbox{|-уайв-|} копать, рыть

\textbf{кудант} \textit{числ.} шестой \mbox{($\leftarrow$\,кудт, -ант)}

\textbf{кӯһт} \textit{числ.} два

\textbf{кӯһтэмплоагкь} \textit{числ.} двенадцать \mbox{($\leftarrow$\,кӯһт, -емп-, лоагкь)}

\textbf{кӯсс} \textit{сущ.} \mbox{|-ӯз-, -ӯссь-|} ель, ёлка

\textbf{куэһцьвэтҍ} \textit{сущ.} \mbox{$\rightarrow$\,куэһцьев, -этҍ}

\textbf{куэһцьев} \textit{сущ.} \mbox{|о.\,ч.|} хвоя; ветка хвойного дерева

\textbf{Куэллнэгк} \textit{собст.} \mbox{|-эг-, -эгкь-|} Куэллнэгк, Кола

\textbf{Куэллнэгке} \textit{собст.} \mbox{$\rightarrow$\,Куэллнэгк, -э}

\textbf{куэссь} \textit{нареч.} когда

\textbf{кыджянт} \textit{числ.} седмой \mbox{($\leftarrow$\,кыджемь, -ант)}

\textbf{кырьйхэй} \textit{сущ.} писатель, пишещий \mbox{($\leftarrow$\,кыррьй, -х-, -эй)}

\textbf{кырьйхэенҍ} \textit{сущ.} \mbox{$\rightarrow$\,кырьйхэй, -энҍ}

\vspace*{\baselineskip}

\textbf{Ладэмыр} \textit{собст.} \mbox{|о.\,ч.|} Ладэмыр, Владимир

\textbf{ле̄в} \textit{глаг.} \mbox{$\rightarrow$\,лӣйе, -эв}

\textbf{лӣ} \textit{глаг.} (\ksdpnumber{3}л. ед.\,ч. наст.вр.) \mbox{$\rightarrow$\,лӣйе}

\textbf{лӣһчемь} \textit{глаг.} \mbox{$\rightarrow$\,лӣйе, -ч-, -эмь}

\textbf{лӣйе} \textit{глаг.} \mbox{|ля̄, лӣ, ля̄йй, ля̄й|} быть

\textbf{лӣйенҍ} \textit{глаг.} \mbox{$\rightarrow$\,лӣйе, -энҍ}

\textbf{лӣййма} \textit{глаг.} \mbox{$\rightarrow$\,лӣйе, -ма}

\textbf{ло̄гкэ} \textit{глаг.} \mbox{|-о̄г-, -оагк-, о̄гкь-|} читать; считать

\textbf{ло̄гкэмь} \textit{глаг.} \mbox{$\rightarrow$\,ло̄гкэ, -эмь}

\textbf{лӯппьтэ} \textit{глаг.} (\ksdpnumber{1}л. ед.\,ч. прош.\,вр.) \mbox{$\rightarrow$\,лӯппьтэ}

\textbf{лӯппьтэ} \textit{глаг.} \mbox{|-ӯпт-, -ӯппт-|} окончиться

\textbf{Луя̄ввьр} \textit{собст.} \mbox{|-я̄вьр-, -я̄ввр-|} Луя̄ввьр, Ловозеро

\textbf{Луя̄вьрэсьт} \textit{собст.} \mbox{$\rightarrow$\,Луя̄ввьр, -эсьт}

\textbf{лыһкма} \textit{глаг.} \mbox{$\rightarrow$\,лыһкэ, -ма}

\textbf{лыһкмэнч} \textit{глаг.} \mbox{$\rightarrow$\,лыһкэ, -ма, -энч}

\textbf{лыһкэ} \textit{глаг.} \mbox{|-ыгк-, -ыһкь-|} работать; делать; трудиться

\textbf{лыһкэв} \textit{глаг.} \mbox{$\rightarrow$\,лыһкэ, -эв}

\textbf{ля̄} \textit{глаг.} (\ksdpnumber{1}л. ед.\,ч. прош.\,вр.) \mbox{$\rightarrow$\,лӣйе}

\textbf{ля̄йй} \textit{глаг.} (\ksdpnumber{3}л. ед.\,ч. прош.\,вр.) \mbox{$\rightarrow$\,лӣйе}

\vspace*{\baselineskip}

\textbf{-ма} \textit{оконч.} прич. 

\textbf{маӈӈа} \textit{нареч.} после, потом, затем

\textbf{мастерскэ} \textit{прилаг.} мастерской

\textbf{ма̄тҍхэ} \textit{глаг.} \mbox{|ма̄тҍха, ма̄тӓһт|} учить, обучать \mbox{($\leftarrow$\,ма̄һтҍ-, -х-, -э)}

\textbf{мӣлльтэ} \textit{послел.} с, за, по

\textbf{мӣн} \textit{местоим.} (\ksdpnumber{1}л. мн.\,ч. ген.) \mbox{$\rightarrow$\,мыйй}

\textbf{мӣнэ} \textit{местоим.} (\ksdpnumber{1}л. мн.\,ч. ген.) \mbox{$\rightarrow$\,мыйй, -э}

\textbf{мӣнэнҍ} \textit{местоим.} (\ksdpnumber{1}л. мн.\,ч. лок.) \mbox{$\rightarrow$\,мыйй, -энҍ}

\textbf{мӣнэтҍ} \textit{местоим.} (\ksdpnumber{1}л. мн.\,ч. акк.) \mbox{$\rightarrow$\,мыйй, -этҍ}

\textbf{моазэ} \textit{нареч.} обратно, назад

\textbf{моаст} \textit{нареч.} обратно, назад

\textbf{мудта} \textit{местоим.} другой, иной

\textbf{мун} \textit{местоим.} (\ksdpnumber{1}л. ед.\,ч. ген.\,/\,акк.) \mbox{$\rightarrow$\,мунн}

\textbf{мунн} \textit{местоим.} я (\ksdpnumber{1}л. ед.\,ч. ном.)

\textbf{Мӯрман} \textit{собст.} Мӯрман, Мурманск

\textbf{Мӯрманскэ} \textit{прилаг.} Мурманский

\textbf{Мӯрмнэ,} \textit{собст.} \mbox{$\rightarrow$\,Мӯрман, -э}

\textbf{Мӯрмнэсьт} \textit{собст.} \mbox{$\rightarrow$\,Мӯрман, -эсьт}

\textbf{мушта} \textit{глаг.} \mbox{$\rightarrow$\,мушшьтэ}

\textbf{мушшьтэ} \textit{глаг.} \mbox{|-ушт-, -ушшт-, -ушьт-|} помнить

\textbf{мыйй} \textit{местоим.} мы (\ksdpnumber{1}л. мн.\,ч. ном.)

\textbf{мэ̄нн} \textit{местоим.} (ген.\,/\,акк. ед.\,ч.) \mbox{$\rightarrow$\,мӣ}

\textbf{мэ̄ннэ} \textit{глаг.} \mbox{|-э̄н-, -э̄нҍ-, -э̄ннҍ-|} идти, пойти

\textbf{мэ̄ннӭ} \textit{глаг.} (\ksdpnumber{1}л. ед.\,ч. прош.\,вр.) \mbox{$\rightarrow$\,мэ̄ннэ}

\textbf{мэ̄нэ} \textit{глаг.} \mbox{$\rightarrow$\,мэ̄ннэ, -э}

\vspace*{\baselineskip}

\textbf{народтҍ} \textit{сущ.} \mbox{|-одҍ-|} народ

\textbf{народӭнч} \textit{сущ.} \mbox{$\rightarrow$\,народтҍ, -энч}

\textbf{-недтӭ} \textit{оконч.} парт. 

\textbf{не̄лльй} \textit{числ.} \mbox{|-е̄льй-, -я̄льй-, -ӣльй-|} четыре

\textbf{ниййт} \textit{сущ.} \mbox{|-ӣйт-|} дочь; девушка

\textbf{нийтнедтӭ} \textit{сущ.} \mbox{$\rightarrow$\,ниййт, -недтӭ}

\textbf{ныдтҍ} \textit{нареч.} да; так

\textbf{нымьп} \textit{числ.} второй; другой

\textbf{нюэзе} \textit{прилаг.} плохо

\vspace*{\baselineskip}

\textbf{о̄һпнувве} \textit{глаг.} (\ksdpnumber{1}л. ед.\,ч. прош.\,вр.) \mbox{$\rightarrow$\,о̄һпнуввэ, -е}

\textbf{о̄һпнуввенҍ} \textit{глаг.} \mbox{$\rightarrow$\,о̄һпнуввэ, -енҍ}

\textbf{о̄һпнуввэ} \textit{глаг.} \mbox{|-ув-, -уввь-|} учиться, обычаться \mbox{($\leftarrow$\,о̄һп-, -н-, -увв-, -э)}

\textbf{о̄һпнуввэв} \textit{глаг.} \mbox{$\rightarrow$\,о̄һпнуввэ, -эв}

\textbf{о̄һпнувэшкӯдтӭ} \textit{глаг.} (\ksdpnumber{1}л. ед.\,ч. прош.\,вр.) \mbox{$\rightarrow$\,о̄һпнуввэ, -шкуэдтҍ-}

\textbf{о̄ллма} \textit{сущ.} \mbox{|о̄ллмэ|} человек

\textbf{о̄ллматҍ} \textit{сущ.} \mbox{$\rightarrow$\,о̄ллма, -этҍ}

\textbf{отделение} \textit{сущ.} \mbox{|о.\,ч.|} отделение

\textbf{отделениясьт} \textit{сущ.} \mbox{$\rightarrow$\,отделение, -эсьт}

\vspace*{\baselineskip}

\textbf{пай} \textit{нареч.} всегда, всё, всё время

\textbf{пайнэмушш} \textit{сущ.} \mbox{|-уж-, -шшь-|} повышение, прогресс \mbox{($\leftarrow$\,пайнэ, -мушш)}

\textbf{пайнэмушшэ} \textit{сущ.} \mbox{$\rightarrow$\,пайнэмушш, -э}

\textbf{парна} \textit{сущ.} дети, детвора (мн.\,ч.) \mbox{($\leftarrow$\,па̄рнэнч)} 

\textbf{па̄рнэнч} \textit{сущ.} \mbox{$\rightarrow$\,па̄ррьн, -энч}

\textbf{па̄ррнатҍ} \textit{сущ.} \mbox{$\rightarrow$\,па̄ррьн, -энч, -этҍ}

\textbf{па̄ррнэтҍ} \textit{сущ.} \mbox{$\rightarrow$\,па̄ррьн, -этҍ}

\textbf{па̄ррьн} \textit{сущ.} \mbox{|-а̄рьн-, -а̄рн-, -арн-|} ребёнок, мальчик, парень

\textbf{па̄ррьшэнч} \textit{сущ.} \mbox{|па̄ррьшя|} ребёнок, мальчик, малыш \mbox{($\leftarrow$\,па̄ррьн, -ш-, -энч)}

\textbf{па̄ррьшэнче} \textit{сущ.} \mbox{$\rightarrow$\,па̄ррьшэнч, -э}

\textbf{па̄ртия} \textit{сущ.} \mbox{|о.\,ч.|} партия

\textbf{педагогическэ} \textit{прилаг.} педагогический

\textbf{пе̄рвэ} \textit{числ.} \mbox{|о.\,ч.|} первый

\textbf{пионерр} \textit{сущ.} \mbox{|-ер-, -еррь-|} пионер

\textbf{пионерэтҍ} \textit{сущ.} пионерр, -этҍ

\textbf{поаһкэ} \textit{глаг.} \mbox{|-оагк-, -оаһкь-|} белеть

\textbf{поаһкэшь} \textit{глаг.} поаһкэ, -эшь

\textbf{подготовительнэ} \textit{прилаг.} подготовительный

\textbf{пра̄ктикум} \textit{сущ.} \mbox{|о.\,ч.|} практикум

\textbf{производственнэ} \textit{прилаг.} производственный

\textbf{пуаз} \textit{сущ.} \mbox{|-ӯдз-|} домашний (северный) олень

\textbf{пӯдзэ} \textit{сущ.} \mbox{$\rightarrow$\,пуаз, -э}

\textbf{пӯдзэтҍ} \textit{сущ.} \mbox{$\rightarrow$\,пуаз, -этҍ}

\textbf{пӯдтэ} \textit{глаг.} \mbox{|-ӯд-, -ӯдтҍ-|} кончиться, иссякнуть

\textbf{пӯдтӭ} \textit{глаг.} (\ksdpnumber{1}л. ед.\,ч. прош.\,вр.) \mbox{$\rightarrow$\,пӯдтэ}

\textbf{пуэдтҍлэ} \textit{глаг.} прибыть, прийти, приехать, поступить \mbox{($\leftarrow$\,пуэдтҍ-, -л-, -э)}

\textbf{пуэдтэль} \textit{глаг.} (\ksdpnumber{3}л. ед.\,ч. прош.\,вр.) \mbox{$\rightarrow$\,пуэдтҍлэ}

\textbf{пуэраст} \textit{прилаг.} хорошо

\textbf{пыдтӭ} \textit{глаг.} \mbox{|-ыд-, -ыдт-|} топить, растапливать

\textbf{пыдтӭмь} \textit{глаг.} \mbox{$\rightarrow$\,пыдтӭ, -эмь}

\textbf{пынне} \textit{глаг.} (\ksdpnumber{1}л. ед.\,ч. прош.\,вр.) \mbox{$\rightarrow$\,пынне}

\textbf{пынне} \textit{глаг.} \mbox{|-ын-, -ынн-, -ынь-|} охранять; беречь; спасать; пасти

\textbf{пэрямп} \textit{нареч.} лучше \mbox{($\leftarrow$\,пэре, -амп)}

\textbf{пэҏт} \textit{сущ.} (ген.\,/\,акк. ед.\,ч., ном мн.\,ч.) \mbox{$\rightarrow$\,пэ̄ҏҏт}

\textbf{пэ̄ҏҏт} \textit{сущ.} \mbox{|-эҏт-, -э̄ҏҏь-|} дом, здание, жилище

\textbf{пя̄дцлушшэ} \textit{глаг.} \mbox{|-уж-, -ушшь-|} заботиться, бескокоиться \mbox{($\leftarrow$\,пя̄дцэль, -л-, -ушш-, -э)}

\textbf{пя̄дцлушшэв} \textit{глаг.} \mbox{$\rightarrow$\,пя̄дцлушшэ, -эв}

\vspace*{\baselineskip}

\textbf{рай} \textit{нареч.} мимо

\textbf{ро̄бхушшэ} \textit{глаг.} \mbox{|-уж-, -ушшь-|} работать \mbox{($\leftarrow$\,ро̄бэһт, -х-, -ушш-, -э)}

\textbf{рӯшш} \textit{сущ.} \mbox{|-ӯш-, -ӯшшь-|} русский

\vspace*{\baselineskip}

\textbf{са̄мь} \textit{сущ.} \mbox{|-а̄ммь-, -оамь-, -а̄мм-|} саами, саамский

\textbf{са̄ррнэ} \textit{глаг.} \mbox{|-а̄рн-, -оаррн-|} говорить; разговаривать

\textbf{се̄йе} \textit{глаг.} \mbox{|о.\,ч.|} сеать, засеять

\textbf{се̄мя} \textit{сущ.} \mbox{|о.\,ч.|} семя

\textbf{се̄мятҍ} \textit{сущ.} \mbox{$\rightarrow$\,семя, -этҍ}

\textbf{сӣнэ} \textit{местоим.} (\ksdpnumber{3}л. мн.\,ч. ген.) \mbox{$\rightarrow$\,сыйй, -э}

\textbf{соаменҍ} \textit{сущ.} \mbox{$\rightarrow$\,са̄мь, -энҍ}

\textbf{соарнэ} \textit{глаг.} (\ksdpnumber{1}л. ед.\,ч. прош.\,вр.) \mbox{$\rightarrow$\,са̄ррнэ}

\textbf{соарнэнҍ} \textit{глаг.} \mbox{$\rightarrow$\,са̄ррнэ, -энҍ}

\textbf{советскэ} \textit{прилаг.} советский

\textbf{со̄нн} \textit{местоим.} \ksdpnumber{1}л. ед.\,ч. ном. (он, она)

\textbf{соннӭ} \textit{местоим.} \mbox{$\rightarrow$\,со̄нн, -э}

\textbf{сруэһк} \textit{сущ.} срок

\textbf{сруэк} \textit{сущ.} (ген.\,/\,акк. ед.\,ч.) \mbox{$\rightarrow$\,сруэһк}

\textbf{сыйй} \textit{местоим.} они

\textbf{сыййт} \textit{сущ.} \mbox{|-ыйт-|} село, селение; деревня

\textbf{сыййта} \textit{сущ.} \mbox{$\rightarrow$\,сыййт, -а}

\textbf{сыйтэ} \textit{сущ.} \mbox{$\rightarrow$\,сыййт, -э}

\textbf{сыйтэсьт} \textit{сущ.} \mbox{$\rightarrow$\,сыййт, -эсьт}

\textbf{-сьт} \textit{оконч.} \mbox{$\rightarrow$\,-эсьт}

\vspace*{\baselineskip}

\textbf{тамьпэ} \textit{нареч.} там

\textbf{танна} \textit{нареч.} тогда

\textbf{темэтҍ} \textit{союз} поэтому

\textbf{тӯййшувв} \textit{глаг.} (\ksdpnumber{3}л. ед.\,ч. наст.\,вр.) \mbox{$\rightarrow$\,тӯййшуввэ}

\textbf{тӯййшуввэ} \textit{глаг.} \mbox{|-ув-, -уввь-|} делаться, работаться \mbox{($\leftarrow$\,тӯйй, -ш-, -увв-, -э)}

\textbf{тыста} \textit{нареч.} здесь

\textbf{тэдта} \textit{местоим.} (ед.\,ч. ном.) этот, эта, это

\textbf{тэнн} \textit{местоим.} (ед.\,ч. ген.\,/\,акк.) этот, эта, это, этого, этой

\textbf{тэсьт} \textit{местоим.} (ед.\,ч. лок.) в этом, в этой, из/от этого, из/от этой

\vspace*{\baselineskip}

\textbf{удць} \textit{прилаг.} маленький

\textbf{училища} \textit{сущ.} \mbox{$\rightarrow$\,училище, -а}

\textbf{учтель} \textit{сущ.} \mbox{|о.\,ч.|} учитель

\textbf{учтельницэнҍ} \textit{сущ.} $\rightarrow$\,учтель, -ница, -энҍ\\

\vspace*{\baselineskip}

\textbf{-ха} \textit{оконч.} абесс. ед.\,ч. 

\vspace*{\baselineskip}

\textbf{це̄ӆӆькэ} \textit{глаг.} (\ksdpnumber{1}л. ед.\,ч. прош.\,вр.) сказать; сообщить; предупредить; назвать

\textbf{цӣӆӆькэ} \textit{глаг.} \mbox{|-я̄ӆк-, -я̄ӆӆк-, -ӣӆӆьк|} \mbox{$\rightarrow$\,це̄ӆӆькэ}

\vspace*{\baselineskip}

\textbf{-ч-} \textit{суфф.} сосл.\,накл. 

\textbf{Чудзья̄ввьр} \textit{собст.} \mbox{|-я̄вьр-, -я̄ввр-|} Чудзья̄ввьр

\textbf{Чудзья̄вьрэсьт} \textit{собст.} \mbox{$\rightarrow$\,Чудзья̄ввьр, -эсьт}

\vspace*{\baselineskip}

\textbf{Ша̄шка} \textit{собст.} \mbox{|о.\,ч.|} Ша̄шка, Александр

\textbf{шко̄ла} \textit{сущ.} \mbox{|о.\,ч.|} школа

\textbf{шко̄лайе} \textit{сущ.} \mbox{$\rightarrow$\,шко̄ла, -йе}

\textbf{шко̄ласьт} \textit{сущ.} \mbox{$\rightarrow$\,шко̄ла, -эсьт}

\textbf{-шкуэдтҍ-} \textit{суфф.} \mbox{|-уад-, уадт-, -ӯдтҍ-|} инк.

\textbf{штобэ} \textit{союз} чтобы

\textbf{штэ} \textit{союз} что

\textbf{шӯрр} \textit{прилаг.} большой

\textbf{шэ} \textit{част.} же

\textbf{шэннтэ} \textit{глаг.} \mbox{|-а̄нт-, -аннт-, -эннҍт-|} расти; рождаться; становиться; выращивать

\textbf{шэнтченҍ} \textit{глаг.} \mbox{$\rightarrow$\,шэннтэ, -ч-, -энҍ}

\vspace*{\baselineskip}

\textbf{ыгкь} \textit{сущ.} \mbox{|ыгь-, ыгк-|} год

\textbf{ыгь} \textit{сущ.} (ген.\,/\,акк. ед.\,ч., ном мн.\,ч.) \mbox{$\rightarrow$\,ыгь}

\textbf{ыжэм} \textit{сущ.} (коми-) ижемец, ижемский

\textbf{Ылльй} \textit{собст.} \mbox{|-ыльй-|} Ылльй, Ильич

\vspace*{\baselineskip}

\textbf{-э} \textit{оконч.} \ksdpnumber{1}л. ед.\,ч. прош.\,вр. 

\textbf{-э} \textit{оконч.} \ksdpnumber{3}л. ед.\,ч. прош.\,вр. 

\textbf{-э} \textit{оконч.} ген. ед.\,ч. 

\textbf{-э} \textit{оконч.} ген. мн.\,ч. 

\textbf{-э} \textit{оконч.} инф. 

\textbf{-э} \textit{оконч.} пов. накл. мн.\,ч. 

\textbf{-эбпе} \textit{оконч.} \ksdpnumber{2}л. мн.\,ч. наст.\,вр. 

\textbf{-эв} \textit{оконч.} \ksdpnumber{3}л. мн.\,ч. наст.\,вр.

\vspace*{\baselineskip}

\textbf{-ӭ} \textit{оконч.} \mbox{$\rightarrow$\,-э}

\textbf{-ӭв} \textit{оконч.} \mbox{$\rightarrow$\,-эв}

\vspace*{\baselineskip}

\textbf{Эвван} \textit{собст.} \mbox{|о.\,ч.|} Эвван, Иван

\textbf{э̄ввтэль} \textit{нареч.} прежде, раньше, сначала

\textbf{эвтас} \textit{нареч.} вперёд \mbox{($\leftarrow$\,э̄ввт, -с)}

\textbf{-эгуэйм} \textit{оконч.} ком. мн.\,ч. 

\textbf{эйй} \textit{глаг.} отриц. \ksdpnumber{3}л. ед.\,ч. 

\textbf{-эк} \textit{оконч.} \ksdpnumber{2}л. ед.\,ч. прош.\,вр. 

\textbf{эле} \textit{глаг.} отриц. повел.\,накл. ед.\,ч. 

\textbf{эллькэ} \textit{глаг.} (\ksdpnumber{1}л. ед.\,ч. прош.\,вр.) \mbox{$\rightarrow$\,аллькэ}

\textbf{-эм} \textit{оконч.} прич. 

\textbf{-эмь} \textit{суфф.} прилаг.

\textbf{-энҍ} \textit{оконч.} \ksdpnumber{3}л. мн.\,ч. прош.\,вр. 

\textbf{-энҍ} \textit{оконч.} есс. 

\textbf{-энҍ} \textit{оконч.} ком. ед.\,ч. 

\textbf{-энҍ} \textit{оконч.} лок. мн.\,ч. 

\textbf{-энч} \textit{суфф.} ласк.\,ф. 

\textbf{-эпь} \textit{оконч.} \ksdpnumber{1}л. мн.\,ч. наст.\,вр. 

\textbf{-эсът} \textit{оконч.} лок. ед.\,ч. 

\textbf{-этҍ} \textit{оконч.} \ksdpnumber{2}л. мн.\,ч. прош.\,вр. 

\textbf{-этҍ} \textit{оконч.} акк. мн.\,ч. 

\textbf{-этҍ} \textit{оконч.} илл. мн.\,ч. 

\textbf{-этҍ} \textit{оконч.} неопр.\,личн.\,ф. наст.\,вр. 

\textbf{этажэнҍ} \textit{прилаг.} этажный

\textbf{-эха} \textit{оконч.} абесс. мн.\,ч. 

\textbf{-эшь} \textit{оконч.} неопр.\,личн.\,ф. прош.\,вр. 

\vspace*{\baselineskip}

\textbf{юррьтэ} \textit{глаг.} (\ksdpnumber{1}л. ед.\,ч. прош.\,вр.) \mbox{$\rightarrow$\,юррьтэ}

\textbf{юррьтэ} \textit{глаг.} \mbox{|-юрт-, -юррт-, -юрьт-|} думать, мыслить

\textbf{юрьтэшкӯдтӭ} \textit{глаг.} \mbox{$\rightarrow$\,юррьтэ, -шкуэдтҍ-}

\vspace*{\baselineskip}

\textbf{я} \textit{союз} и

\textbf{-я} \textit{оконч.} \mbox{$\rightarrow$\,-а}

\textbf{-я} \textit{суфф.} \mbox{$\rightarrow$\,-а}

\textbf{-яһта} \textit{оконч.} \mbox{$\rightarrow$\,-ха}

\textbf{я̄ӆчемь} \textit{глаг.} \mbox{$\rightarrow$\,е̄лле, -ч-, -эмь}

\textbf{я̄ӆченҍ} \textit{глаг.} \mbox{$\rightarrow$\,е̄лле, -ч-, -энҍ}

\textbf{я̄мме} \textit{глаг.} \mbox{|я̄м-, я̄мм-, ёам-|} умирать, умереть; неметь, онеметь

\end{multicols}

\newpage

\noindent Сокращений\\

\noindent\textit{\ksdpnumber{1}л.} — первое лицо\\
\textit{\ksdpnumber{2}л.} — второе лицо\\
\textit{\ksdpnumber{3}л.} — третое лицо\\
\textit{абесс.} — абессив	(лишительный падеж)\\
\textit{акк.} — аккузатив	(винительный падеж)\\
\textit{ген.} — генитив	(родительный падеж)\\
\textit{глаг.} — глагол\\
\textit{ед.\,ч.} — единственное число\\
\textit{есс.} — ессив	(превратительный падеж)\\
\textit{илл.} — иллатив	(дательно-направительный падеж)\\
\textit{инк.} — инкоатив (начинательный глагол)\\
\textit{инф.} — инфинитив \\
\textit{ком.} — комитатив	(совместный падеж)\\
\textit{ласк.\,ф.} — ласкательная форма (диминутив)\\
\textit{лок.} — локатив	(местный падеж)\\
\textit{мест.} — местоимение\\
\textit{мн.\,ч.} — множественное число\\
\textit{нареч.} — наречие\\
\textit{наст.\,вр.} — настоящее время\\
\textit{неопр.\,личн.\,ф.} — неопределенно-личная форма\\
\textit{о.\,ч.} — отсутствие чередования\\
\textit{оконч.} — окончание\\
\textit{отриц.} — отрицательная форма\\
\textit{парт.} — партитив\\
\textit{повел.\,накл.} — повелительное наклонение (императив)\\
\textit{послел.} — послелог\\
\textit{прилаг.} — прилагательное\\
\textit{прич.} — причастие\\
\textit{прош.\,вр.} — прошедшее время\\
\textit{собст.} — имя собственное \\
\textit{сосл.\,накл.} — сослагательное наклонение (конъюунктив)	\\
\textit{суфф.} — суффикс\\
\textit{сущ.} — существительное\\
\textit{числ.} — имя числительное

\newpage
\fancyhead{}
\fancyhead[RO]{\thepage}

\clearpage
\hbox{}
\thispagestyle{plain}

\clearpage
\hbox{}
\thispagestyle{empty}

\end{document}